% PM 트랙 Day 3 슬라이드 본문 — 손으로 쓴 정본(한 프레임 = 한 쪽). Day2_슬라이드_body.tex 와 같은 형식.
% 래퍼: Day3_슬라이드.tex(슬라이드판) / Day3_슬라이드_노트.tex(노트판, 우측에 \note).
% 화면에는 결론만, 설명·근거·예외는 각 프레임의 \note{}에. Day3_슬라이드.md 는 이 파일의 텍스트 미러다.
% 그림은 ../그림/PM_Day3/*.pdf (벡터). 그림 번호 = 파일명 번호(그림 2-3 = fig_2_3_*). 모든 숫자는 Day3 워크북 완성 예시본(정본)에서.
% 데이터 일자 2026-09-30. BAC 145.80 / 집행 한도 156.00 / PV 77.83 / EV 73.43 / AC 71.13 / SPI 0.943 / SPI(t) 0.920 / CPI 1.032(조정 0.968) / EAC4 150.86.


\newcolumntype{L}[1]{>{\raggedright\arraybackslash}p{#1}}
\newcommand{\ra}{$\rightarrow$}
\newcommand{\til}{\textasciitilde}
\newcommand{\keymsg}[1]{\par\vfill{\color{mDarkTeal}\bfseries #1}\par}
\newcommand{\figcap}[1]{\par\smallskip{\footnotesize\color{black!60}#1}\par}
\newcommand{\fig}[2][0.66]{\centering\includegraphics[height=#1\textheight,width=\textwidth,keepaspectratio]{#2}\par}
\newcommand{\subhead}[1]{{\color{mDarkTeal}\bfseries #1}\par}
\newenvironment{tbl}[2][\small]{\begin{center}\usebeamercolor[fg]{normal text}\setlength{\tabcolsep}{4pt}\renewcommand{\arraystretch}{1.15}#1\begin{tabular}{#2}\toprule}{\bottomrule\end{tabular}\end{center}}

% =====================================================================
% 도입

% =====================================================================
% 도입
% =====================================================================
\begin{frame}{PM 트랙 Day 3 — 어제 세운 자(尺)에 오늘 실적을 댄다}
\begin{itemize}
\item 프로젝트 매니지먼트 4일 특강 · 3일차 · 실행과 통제(Executing · Monitoring \& Controlling)
\item 사례는 같다 — 온담F\&B홀딩스 ONE ONDAM \textbf{P1} · 승인 168.0 / 집행 한도 156.0 / PMB 145.8억 · FP 12,400 고정가
\item 시점 고정 — \textbf{데이터 일자 2026-09-30}. 착수 9개월, S19 진행 중(R2 완료 직후), 규제 ① 검증 통과 직후, G2(11-27) 앞
\item 어제 등록부의 리스크가 실제로 일어났다 — R-05 부분 실현(구현 3주 지연), R-04 대응(라이선스 50/50), R-03 트리거 발동(폴백 확정)
\item Day2 다섯 기준선(⑥\til{}⑩)이 오늘 다섯 통제 산출물(⑪\til{}⑮)의 \textbf{자}다
\end{itemize}
\keymsg{오늘의 한 문장 — 통제는 편차를 재는 일이 아니라, 편차에 대해 누가 무엇을 결정해야 하는지를 적는 일이다}
\note{표지. 어제 워크북 완성 예시본 5종 — 특히 §3.5 CPM 계산표·스프린트 캘린더, §4.5 월별 PV 표(2026-09 누적 77.83), §5.5 등록부(잔여 EMV 6.78)와 격자 7축 — 를 읽어 왔는지 확인한다. 읽지 않은 수강생은 워크북 §0.3 한 페이지 요약(2026-09-30까지 일어난 일 아홉 줄)과 §0.4 "Day2에서 가져오는 것" 표를 4교시 전까지 읽도록 안내한다. Day1이 "결정되지 않은 것을 결정하는 문서", Day2가 "결정된 것을 분해하고 숫자로 만드는 문서"였다면, Day3는 "기준선과 실적의 차이를 재고 그 차이에 대해 누가 언제까지 무엇을 결정해야 하는지를 적는 문서"다. 그래서 오늘의 모든 산출물은 결정 요청 칸으로 끝난다. 시나리오는 [추가 설정]이며 숫자는 Day2 기준선과 정합하게 정했다 — 워크북과 교재가 다르면 워크북이 정본. 예상 소요 3분.}
\end{frame}

\begin{frame}[shrink=10]{Day2 기준선 5종 \ra{} Day3 통제 산출물 5종}
\fig[0.62]{fig_0_1_day2_to_day3}
\figcap{그림 0-1. 성과보고(⑫)가 중심이고 나머지 넷이 그 각주다. 되돌아가는 파선은 승인된 변경(BL-02)으로만 기준선이 움직였음을 보인다 — 지연(A10 $-$10)과 초과(밀도 0.132)는 편차로 남아 있다}
\keymsg{Day1 §0 마스터 표의 Day3 3종이 Day2 §7의 요구(품질 통제)를 받아 5종으로 확정되었다 — 4일 20종}
\note{교재 서문의 표와 그림 0-1. ⑪ 변경요청서·영향 분석·CCB 의결서는 ⑥ WBS 사전 카드(추정·의존·인수 기준)·⑦ RTM 변경 이력 열·⑩ 격자(±35건·FP ±372·원가 10.2·리스크 8억)와 Day2 §6의 헌장 v1.1 수정 요청 10건 중 1·2·4번을 받는다. ⑫ 성과보고는 ⑨ 월별 PV 표 17개월(BAC = PMB 145.8)과 ⑧ CPM 계산표·스프린트 캘린더를. ⑬ 이슈 로그·리스크 등록부 갱신은 ⑩ 등록부(잔여 6.78·트리거·2차 리스크 예비 번호 R-14\til{}R-18)와 컨틴전시 배정표(4.60 + 5.60)를. ⑭ 품질 통제 기록은 ⑦ 시험·검수 열과 ⑩ 품질 허용범위 두 계층(작업패키지 0.13 / 프로젝트 0.4 + 심각도 1·2 미해결 0)을. ⑮ 조달·계약 변경 기록은 ⑨ 단가 55만 원/FP·⑥ L2 FP 배분·⑧ 여유 소유권 문안을. 세 가지를 미리 말한다 — 기준선 값을 바꾸지 말고 인용하라(틀렸음을 발견하면 "측정 규칙 차이"로 분리해 적고 CR-10으로 CCB에) / 오늘의 시나리오는 나쁜 소식을 담고 있다(좋은 소식만 담은 사례는 연습이 안 된다) / 모든 문서는 결정 요청으로 끝난다. 예상 소요 4분.}
\end{frame}

\begin{frame}{오늘의 흐름}
\begin{columns}[T]
\begin{column}{0.47\textwidth}
\subhead{오전 — 이론 (제1\til{}3장)}
\begin{itemize}
\item \textbf{1교시} 90분 · 통제론 · 성과 측정 EVM — PMB·BAC·Rules of Credit·편차와 지수 — 제1장, 제2장 2.1\til{}2.3
\item \textbf{2교시} 90분 · EAC 4공식 · Earned Schedule · 착시 5개 · 흐름 지표 · 임계경로 $-$10과 EX-01 — 제2장 2.3\til{}2.9
\item \textbf{3교시} 90분 · 변경 관리 — 6축 · 처리 경로 · BL-02 vs 재기준선 · 변경 대가 — 제3장
\end{itemize}
\end{column}
\begin{column}{0.51\textwidth}
\subhead{오후 — 실습과 통제의 나머지}
\begin{itemize}
\item \textbf{4교시} 90분 · \textbf{실습 ①} 변경요청서·영향 분석·CCB + 성과보고 EVM 손계산 — 워크북 §1\til{}2
\item \textbf{5교시} 75분 · 품질 · 이슈·리스크 · 조달 통제 — 결함 곡선·감리 귀속·EMV 재평가·검수 보류 회신·귀책 분해 — 제4\til{}5장
\item \textbf{6교시} 75분 · \textbf{실습 ②} 이슈 로그·품질 기록·감리 대응·검수 보류 회신 — 워크북 §3\til{}5
\item \textbf{마무리} 30분 · 학술 배경 · 읽을 자료 · Day4 예고 — 제6장
\end{itemize}
\end{column}
\end{columns}
\keymsg{오전 세 교시가 성과보고와 변경요청서의 이론이고, 그 둘을 4교시에 손으로 쓴다 — 나머지 셋은 오후에}
\note{Day2와 같은 골격 — 오전 이론 셋, 점심은 3교시 뒤, 오후는 실습 ① \ra{} 이론(품질·이슈·조달) \ra{} 실습 ②. 1교시가 제2장 전반(2.1\til{}2.3 편차·지수까지)을 포함하는 이유는 EVM의 정의와 Rules of Credit이 오전 내내 쓰이기 때문이고, 2교시는 EAC 4공식부터 시작한다. 4교시 실습 ①이 ⑪·⑫ 둘을 맡는 것은 두 문서가 서로의 입력이기 때문 — 변경 4건의 FP 델타(+185)가 EAC4의 ETC 행에 들어가고, 성과보고의 "미반영 변경 CR-03·07"이 변경 로그를 가리킨다. 5교시 75분에 제4장(품질·이슈·리스크·예외 보고)과 제5장(보고 체계·mum effect·조달·계약)을 다 다루므로 속도가 빠르다 — 6교시 실습 ②가 그 내용을 손으로 다시 밟는다. 예상 소요 3분.}
\end{frame}

\begin{frame}{학습 목표 — 오늘이 끝나면}
\begin{columns}[T]
\begin{column}{0.5\textwidth}
\subhead{오전 이론}
\begin{enumerate}
\item 통제 사이클 5단계와 조치 3층을 \textbf{허용범위 격자}로 가른다
\item BAC 145.8을 ANSI/PMI 19-006 용어로 설명한다
\item EAC 4공식·\textbf{SPI(t)·CPM}으로 "SPI가 오르는데 예외 보고"를 설명한다
\item 착시 다섯에 이름을 붙이고 \textbf{병기}한다
\item 변경을 6축으로 재고 CCB 여부를 판정한다
\end{enumerate}
\end{column}
\begin{column}{0.48\textwidth}
\subhead{오후}
\begin{enumerate}\setcounter{enumi}{5}
\item 결함 곡선의 평탄을 시험 노력으로 판독한다
\item 감리 지적의 \textbf{귀속}에 문서로 이의한다
\item 잔여 EMV 변동을 분해하고 컨틴전시 소진을 읽는다
\item 구두 검수 보류에 서면으로 답하고 귀책을 \textbf{지금} 분해한다
\end{enumerate}
\end{column}
\end{columns}
\keymsg{정직하게 적을 셋 — SPI가 오르는데 EX-01 / CPI 1.032가 좋은 소식이 아닌 이유 / 소진 13.8\%가 안심이 아닌 이유}
\note{아홉 목표 중 1\til{}5는 오전(1\til{}3교시), 6\til{}8은 5교시, 9는 두 실습. 오늘의 결론을 미리 말해 둔다(교재 "Day3을 마치며") — 실행 단계에서 정직하게 적어야 할 세 가지. SPI가 0.913 \ra{} 0.943으로 오르는데 예외 보고를 낸 이유는 SPI(t)가 0.920으로 정체이고 임계경로 A10이 $-$10영업일이며 컷오버 창은 허용범위 0이라는 사실. CPI 1.032가 좋은 소식이 아닌 이유는 라이선스 4.70의 인식 기준 차이를 걷어 내면 0.968이고 EAC3$'$ 155.04는 집행 한도까지 0.96 남았다는 사실. 컨틴전시 소진 13.8\%가 안심의 근거가 아닌 이유는 미배정 풀이 스프린트 4회분에서 1.75회분이 되었고 관리예비비 대상 2.20의 원인이 전부 프로그램·스폰서 계층이라는 사실. 이 셋은 나쁜 소식이 아니라 통제가 한 일이며, 스폰서에게 결정을 요청하는 근거다. 그리고 재기준선은 0회다. 예상 소요 3분.}
\end{frame}

% =====================================================================
\section{1교시 (90분) — 통제론 · 성과 측정 EVM: PMB · BAC · Rules of Credit · 편차와 지수}
% =====================================================================

\begin{frame}{이 교시에서 배우는 것}
\begin{itemize}
\item 통제 사이클 다섯 단계 — 측정 \ra{} 분석 \ra{} 예측 \ra{} 조치 \ra{} 기준선 갱신, 그리고 조치의 세 층 [§1.1]
\item 통제의 단위 — 예측형의 게이트, 적응형의 반복, 하이브리드의 세 층(스프린트 / 월 / 게이트)과 층 사이의 연결 규칙 [§1.2]
\item 표준 — PMBOK 8판 두 Focus Area · PRINCE2 7 세 프로세스와 예외에 의한 관리 · ISO 21502 §7 · 한국 SI 네 주기 [§1.3\til{}1.4]
\item 가장 비싼 오류 셋 — 90\% 증후군 · 기준선 없는 보고 · 계획을 고쳐 편차를 없애기 [§1.5]
\item EVM의 전제 PMB와 BAC(145.8 vs 156.0) · Rules of Credit 여섯 기법 · SV·CV·SPI·CPI 읽기 규칙 셋 [§2.1\til{}2.3]
\end{itemize}
\keymsg{자가 있어야 편차가 있고, 편차가 있어야 통제가 있다 — 그러나 편차를 재는 것만으로는 아무 일도 일어나지 않는다}
\note{제1장 전체와 제2장 2.1\til{}2.3 편차·지수까지를 90분에 다룬다. 시간 배분은 §1.1 12분, §1.2 10분, §1.3\til{}1.4 12분, §1.5 8분, §2.1 15분, §2.2 15분, §2.3 편차·지수 10분, 정리 5분(EAC 4공식은 2교시 첫 주제). §2.1 BAC 논쟁과 §2.2 Rules of Credit은 4교시 실습 ① 성과보고 손계산의 직접 입력이므로 여기서 숫자를 정확히 남긴다 — PV 77.83 / EV 73.43 / AC 71.13. 수강생이 PMP 소지자라면 ANSI/PMI 19-006-2019가 2011년 Practice Standard 2판을 대체한 현행 정본임을 판본명으로 정확히 부르는 것이 신뢰의 첫 단계. 예상 소요 2분.}
\end{frame}

\begin{frame}[shrink=15]{통제 사이클 — 측정 \ra{} 분석 \ra{} 예측 \ra{} 조치 \ra{} 기준선 갱신}
\fig[0.66]{fig_1_1_control_cycle}
\figcap{그림 1-1. 다섯 단계와 온담 2026-09-30의 실제 값. 조치의 세 층(PM 전결 / 예외 보고 / 변경요청)을 가르는 것이 Day1의 허용범위 격자, 기준선 갱신은 승인된 변경으로만. 바깥 동심원 = 통제의 주기}
\keymsg{"SV $-$4.40억"은 사실이지 결정이 아니다 — 통제는 편차에 대해 무엇을 할지를 정하고 실행하는 일이다}
\note{§1.1 개념. 측정 — 데이터 일자를 고정하고 기준선과 같은 규칙으로 잰다(인수 FP 7,020 · AC 71.13 · A10 착수 09-28 · 결함 1,140/1,024). 분석 — 크기가 아니라 원인: SV $-$4.40 중 $-$1.96은 측정 시차이고 $-$2.44가 실질 지연이며 그것은 SI 열($-$5.31)에 몰려 있고 시험·감리(+0.70)가 일부 가린다. 예측 — 전제를 적는다: EAC1 141.23은 "지금 효율이 끝까지", EAC4 150.86은 항목별 잔여 근거의 합. 조치 — 세 층: 허용범위 안 PM 전결(결함 해소 전담 2명), 초과 예상 시 예외 보고(EX-01 컷오버 창 이탈), 기준선 변경은 CR \ra{} CCB. 기준선 갱신 — CCB 제10차(10-15)가 CR-08\til{}11을 의결하고 BL-02(10-23)가 발행되면 FP 12,400 \ra{} 12,585, PMB 145.80 \ra{} 146.82; 그러나 A10 $-$10은 기준선에 들어가지 않는다. 주기의 길이(월 성과보고·격주 변경협의체·2주 리뷰·주간 PM 회의)가 통제의 해상도. 예상 소요 6분.}
\end{frame}

\begin{frame}{온담으로 읽기 — 사이클은 아홉 번 돌았고, 두 번 실패하고 한 번 성공했다}
\begin{itemize}
\item \textbf{실패 1 — 4월 제4호}: SPI 0.784를 "주의"로 적고 끝. 열별 분해가 없어 상용SW $-$6.0(조달 이연)과 SI 실질 0이 한 숫자로 — 측정은 되었으나 분석이 없었고, 그래서 조치도 없었다
\item \textbf{실패 2 — 5월 제5호}: CV가 $-$0.16 \ra{} +2.70. "원가 절감"으로 읽혔다. 실제는 03-13 라이선스 50/50 계약 변경의 흔적 — EV 100\% 인식, AC 50\%만 발생. 8월에야 발견
\item \textbf{성공 — 9월 제9호}: S18 3주 연장(09-04) \ra{} A10 착수 09-28 \ra{} TF $-$10 \ra{} as-is MG3 2027-03-05 \ra{} 컷오버 1차 창 이탈 예측 \ra{} \textbf{EX-01 10-02 제출}(인지 09-28 + 4영업일, 규정 5영업일 안)
\item 두 실패와 한 성공의 차이는 도구가 아니라 \textbf{양식}이다 — 4월에 열별 표가, 5월에 조정 행이 있었다면
\item 그래서 제9호는 조기 신호 5개를 사후에 적어 다음 주기의 양식을 바꿨다 — 조정 행 · 시험 노력 열 · "임계경로 TF 영향" 칸
\end{itemize}
\keymsg{통제 사이클의 실무적 정체는 보고서 양식이다 — 양식에 없는 칸은 관측되지 않는다}
\note{§1.1 "온담의 사례로 읽기". 4월의 실패에서 R-04 이슈 전환(IS-01)은 2월에 이미 있었지만 그것이 4월 SPI의 원인이라는 연결이 보고서에 없었다 — 분석 열이 없으면 이슈 로그와 성과보고는 서로를 설명하지 못한다. 5월의 실패는 측정 규칙(AC 인식 기준)이 실행 중에 바뀌었는데 성과보고 양식에 "조정 행"이 없었던 문제이며, 이 착시는 나중에 2027-01 자금 봉우리에 얹혀 R-17(2027 자금 집중)의 일부가 된다 — 착시는 지표만 왜곡하는 것이 아니라 리스크를 만든다. 9월의 성공은 스프린트 캘린더와 CPM 계산표가 같은 날짜를 가리켰기 때문에 가능했다(Day2 3.1절의 세 층). 다음 프레임에서 그 세 층을 본다. 예상 소요 5분.}
\end{frame}

\begin{frame}[shrink=8]{통제의 단위 — 예측형의 게이트, 적응형의 반복, 하이브리드의 세 층}
\fig[0.6]{fig_1_2_three_control_layers}
\figcap{그림 1-2. 같은 사건(S18 3주 연장)이 세 층을 관통하며 다른 해상도로 읽힌다 — 스프린트 층 "이월 420 FP" / 월 층 "A10 TF $-$10, SPI(t) 0.920" / 게이트 층 "컷오버 1차 창 이탈 예상 \ra{} 예외 보고"}
\keymsg{층 사이의 연결 규칙이 없으면 아래층의 작은 결정이 위층의 큰 사건이 될 때까지 아무도 보지 못한다}
\note{§1.2 개념. 그림은 왼쪽에서 오른쪽으로 읽는다 — 맨 아래 스프린트 층(2주 × 36회: 스프린트 리뷰에서 인수 FP 확인, throughput·WIP·work item age, 결정은 이재현 PM과 P1 PMO), 가운데 월 층(성과보고 EVM·ES·CPM TF, 변경협의체 격주·CCB 월, 등록부 갱신, 요구사항 동결 매월 마지막 수요일 — 결정은 P1 PM 전결과 CCB), 맨 위 게이트 층(G0\til{}G5: 재기준선 여부, 컷오버 go/no-go, 헌장 확정 — 스폰서·프로그램 이사회). 예측형의 단위는 게이트와 보고 주기이고 EVM은 그 자연스러운 도구(PMB가 있어야 PV가 있다). 적응형의 단위는 반복과 리뷰이며 "기준선 대비 편차" 개념이 약하고, EVM을 그대로 쓰면 BAC가 흔들리고 스토리포인트는 EV 단위가 못 된다(2.7절). 하이브리드는 둘을 층으로 겹친다. 온담 사례 — 05-12 POS 화면 정의서 승인 대기 시작(R-05 실현): 스프린트 층은 S10을 3주로 늘려 흡수, 월 층은 SI SV $-$2.49를 "약간의 지연"으로 넘겼고 대체 결정권자 지정은 06-03(3주 늦음), 변경협의체는 07-15에 12영업일 중 7영업일 발주사 귀책·0.42억 판정, 게이트 층은 G2 앞 R-05 재평가로만 본다. 예상 소요 6분.}
\end{frame}

\begin{frame}{세 층이 부딪히는 지점 — 연결 규칙 셋}
\begin{itemize}
\item \textbf{스프린트 층은 '인수 FP'로, 월 층은 'EV'로 성과를 말한다} \ra{} 잇는 것은 Rules of Credit(인수 FP ÷ 620 × 1.96). 규칙이 없으면 리뷰의 "완료"와 보고의 "가치"가 다른 것을 가리킨다
\item \textbf{스프린트 층은 범위를 재조정하고 월 층은 범위를 동결한다} \ra{} 잇는 것은 접수 마감(동결 2주 전)과 변경협의체. 없으면 백로그 조정이 범위 변경의 뒷문
\item \textbf{스프린트 층의 지연은 다음 스프린트로 이월되지만 월 층의 지연은 임계경로 TF를 먹는다} \ra{} 잇는 것은 캘린더와 CPM이 같은 날짜를 가리키는 것
\item 그래서 S18 연장 결정(09-04)에 "임계경로 TF 영향" 칸이 있어야 했다 — 있었다면 예외 보고는 09-11에 나갔다(3주 빠름)
\item 조기 신호 2 — 5월 2주차 work item age 12영업일은 스프린트 층 지표가 월 층 대응(대체 결정권자)을 3주 당길 수 있었다는 기록
\end{itemize}
\keymsg{한 층만 있는 조직 — 스프린트 리뷰만 있어 임계경로가 안 보이거나, 월 보고만 있어 지연을 한 달 늦게 안다}
\note{§1.2 후반과 "실무에서 어떻게 틀리는가". 실무 오류 셋: 한 층만 있다 / 층 사이의 연결 규칙이 없다(스프린트 "완료"를 EV로 그대로 써서 착수 시 100\% 인정 — PV와 같아지게, 또는 백로그 조정을 변경 절차 밖에 둔다) / 아래층의 결정에 위층 영향 칸이 없다(스프린트 연장·항목 이월·시험 인력 이동은 일상 결정이지만 그 결정마다 임계경로 TF 영향이 계산되지 않으면 위층은 그 합을 예외 보고 시점에 처음 본다). 온담에서 스프린트 층의 "S10을 3주로" 결정이 월 층의 "A05 TF $-$5"를 만들고, 넉 달 뒤 S18의 $-$10과 합쳐져 게이트 층의 예외 보고가 되었다. 세 층은 각각 제 일을 했다 — 문제는 층 사이였다. 예상 소요 4분.}
\end{frame}

\begin{frame}{표준에서는 — PMBOK 8판 두 Focus Area · 다섯 성과영역 · ISO 21502 §7}
\begin{tbl}[\scriptsize]{L{0.17\textwidth}L{0.58\textwidth}L{0.15\textwidth}}
\textbf{8판 성과영역} & \textbf{실행·통제에서 다루는 것} & \textbf{교재} \\ \midrule
Governance & CCB·보고·에스컬레이션·예외 — 승인된 변경이 기준선과 PMB를 갱신 & 제3장 \\
Schedule & CPM 갱신·TF 소진·마일스톤 전망 + 적응형 지표(번다운·처리량) 병기 & 2.8절 \\
Finance & EVM(EV·PV·AC \ra{} EAC·ETC·TCPI·VAC), 예측, Earned Schedule & 제2장 \\
Risk & 트리거 관측 \ra{} 대응 실행 \ra{} 잔여·2차 평가 \ra{} 예비비 대조 & 4.5절 \\
Stakeholders & 이슈 관리·에스컬레이션·참여 감시 & 4.4·5.3절 \\
\end{tbl}
\begin{itemize}
\item ISO 21502 §7 — 변경 7.10 · 리스크 7.8 · \textbf{이슈 7.9(별도 조항)} · 보고 7.12 · 조달 7.17 — 계약이 중립적으로 인용하는 번호
\end{itemize}
\keymsg{사내 표준이 6판이면 6판 번호(4.5/4.6/6.6/7.4/8.3/11.6/12.3/13.4)로 인용하고 8판 명칭을 병기한다 — 워크북 각 X.2 표}
\note{§1.3 PMBOK 8판·ISO 부분. 8판(2025)은 7판의 원칙·성과영역 위에 프로세스 기반 Focus Area를 되돌려 놓았고, 실행·통제에 해당하는 것은 Executing과 Monitoring and Controlling 두 Focus Area다. 6판의 8.3 Control Quality는 "QC = 산출물 검사, QA = 프로세스 검사"의 이분법이었고 7판이 그것을 지웠으며 8판이 다시 구분하되 인수와의 연결을 앞세운 것 — 온담에서 그 연결은 "심각도 1·2 미해결 0은 G3 인수 조건"이다(제4장). ISO 21502가 이슈 관리(7.9)를 리스크 관리(7.8)와 별도 조항으로 둔 것이 이 교재가 이슈 로그와 리스크 등록부를 한 문서의 두 면으로 다루되 분리하는 근거(4.4절). 실무 오류 — 6판 ITTO를 8판이라고 부르는 것. 예상 소요 5분.}
\end{frame}

\begin{frame}{PRINCE2 7 — 세 프로세스와 예외에 의한 관리(management by exception)}
\fig[0.56]{fig_1_3_management_by_exception}
\figcap{그림 1-3. 허용범위는 아래로 내려오는 위임이고 예외 보고는 위로 올라가는 반대급부다. 세 보고서(Checkpoint·Highlight·Exception)는 주기와 수신자가 다르며, 합치면 예외가 월례에 묻힌다}
\keymsg{예외 보고는 권한의 반납이 아니다 — "이 밖의 결정은 당신의 것이니 대안을 들고 가겠다"는 위임의 반대급부다}
\note{§1.3 PRINCE2 부분. 세 프로세스 — Controlling a Stage(PM의 일상: 작업 패키지 승인, 이슈·리스크 접수, Highlight Report, 시정 조치, 에스컬레이션), Managing Product Delivery(팀 매니저 = 수행사 PM 이재현: 작업 패키지 수락·실행·Checkpoint Report·인도), Managing a Stage Boundary(게이트 앞: 다음 단계 계획, BC 갱신, End Stage Report, 예외 계획). 분리가 주는 것 둘 — 발주사 PM과 수행사 PM의 역할이 문서 수준에서 갈라진다(Checkpoint는 이재현 \ra{} P1 PM, Highlight는 P1 PM \ra{} 정미란), 게이트가 단계를 닫고 여는 결정 지점이라는 것. Progress practice의 관리 방식 — 각 계층은 위가 준 허용범위 안에서 자율 결정하고, 넘을 것으로 예상되는 순간 예외 보고. 허용범위 없는 조직은 모든 결정이 위로 올라가 PM이 없고, 예외 보고 없는 조직은 PM이 권한 밖 결정을 혼자 내려 스폰서가 없다. Exception Report 구조(원인·6개 성과 목표 영향·대안 3·권고·결정 시한)가 헌장 §12 양식이고 EX-01(10-02)이 첫 실물. 실무 오류 — 예외 보고를 "실패 보고"로 이해해 내지 않는 것(늦을수록 대안이 줄고, 대안 없는 예외 보고는 사후 통보), Highlight를 매달 전부 G로 내는 것, Checkpoint와 Highlight를 합쳐 수행사 작업 패키지 상태가 스폰서에게 그대로 가는 것. 예상 소요 6분.}
\end{frame}

\begin{frame}{한국 SI의 통제 사이클 — 네 주기가 세 조직에 걸쳐 있다 · 분야별 대조}
\begin{itemize}
\item \textbf{주간·월간 보고}(PMO) · \textbf{감리}(제3자, 온담 4회 4.2억) · \textbf{검수}(계약·재경) · \textbf{스프린트 리뷰}(수행사) — 네 주기, 세 조직
\item 감리 지적 \ra{} 검수 보류 \ra{} 대기 비용 청구 \ra{} 귀책 분쟁 — 한 사슬. "공정률"은 EV가 아니라 시간 비율인 경우가 대부분
\end{itemize}
\begin{tbl}[\scriptsize]{L{0.12\textwidth}L{0.2\textwidth}L{0.3\textwidth}L{0.28\textwidth}}
\textbf{분야} & \textbf{통제 단위} & \textbf{성과 측정} & \textbf{인수·검수} \\ \midrule
건설·EPC & 기성(progress payment) & 물리적 진도 규칙이 계약 문서, EVM은 계약 요건 & ITP hold point · NCR · 상주 감리 \\
SI(한국) & 릴리스·단계 검수 & 공정률(시간 비율) — FP 인수가 자연스러운 EV 단위 & 감리 지적 \ra{} 검수 \ra{} 기성 \\
제조 NPD & 게이트(DR·PV·양산 이관) & 마일스톤 가중 EV, 양산 원가 전망 & DV·PV 결함 · ECN \\
제약 임상 & 실사·규제 제출 & 등록 환자·사이트 수 — 원가보다 일정 & GCP 실사 · 규제 질의 \\
\end{tbl}
\keymsg{네 분야가 공유하는 것 — 인수의 조건이 통제 지표를 결정한다. 지표는 "돈이 움직이는 조건"이다}
\note{§1.4. 그림 1-4(한국 SI 네 주기와 세 조직, 09-25 감리 통보 \ra{} 10-16 검수 통보 연쇄)는 교재에 있으며 여기서는 표로 대신한다. 온담이 Day2에서 월별 PV 표를 만든 것은 "EVM을 하기 위해서"가 아니라 4월의 SPI 하락이 조달인지 개발인지를 열별로 구분하기 위해서였고(Day2 §4.6 해설), 그 예고가 실제로 일어났다. 실무 오류 — 공정률을 EV로 쓰는 것 / 감리와 검수를 다른 조직의 일로 두고 PM이 연결을 관리하지 않아 감리 보고서 귀속 칸이 검수 분쟁의 유일한 기록이 되는 것 / 주간 보고를 수행사가 쓰고 발주사가 읽기만 하는 것 — 발주사의 결정 SLA 준수율이 어느 문서에도 없으면 변경협의체에서 "발주사가 5월에 12일 미뤘다"에 반박할 자리가 없다(5.4절). 예상 소요 5분.}
\end{frame}

\begin{frame}{가장 흔하고 비싼 오류 셋 — 90\% 증후군 · 기준선 없는 보고 · 계획을 고쳐 편차를 없애기}
\fig[0.52]{fig_1_5_rebaseline_trap}
\figcap{그림 1-5. 왼쪽: BL-01 고정 + BL-02는 승인 변경(+1.02·+185 FP)만 계단으로 — A10 $-$10은 편차로 남는다. 오른쪽(가상): 분기 재기준선 — 마지막 기준선 대비 편차는 0이지만 최초 대비 누적 지연은 어느 문서에도 없다}
\keymsg{셋의 공통점 — 측정 규칙·기준선 버전·편차 이력이 문서에 없다. 통제는 분석 능력이 아니라 기록의 문제이고, 기록은 양식의 문제다}
\note{§1.5. 90\% 증후군 — 담당자 진척률은 종료 3개월 전까지 항상 80\til{}90\%. 거짓말이 아니라 측정 규칙의 부재("진척률"이 활동 수·기간·원가·FP·시험 통과율 중 무엇의 비율인지 정의되지 않음). Kerzner의 진척 착시(6.8절). 처방은 사람이 아니라 Rules of Credit(2.2절) — 진행 중 0(0/100), 물리 단위(인수 FP), 마일스톤 가중. 온담 SI 구현 EV가 "진행 중 스프린트 0"인 이유이며 그것이 만드는 구조적 SV $-$1.96은 착시가 아니라 규칙의 비용. 기준선 없는 보고 — "지난달보다 좋아졌다"는 추세이지 편차가 아니다(방향만 말하고 도착지를 말하지 않음); 버전 미명시(BL-01인가 02인가)는 EV와 PV가 다른 범위를 잴 수 있으니 제9호는 첫 표에 "비교 기준선 BL-01, BL-02 10-23 예정, 미반영 변경 CR-03·07 별도"를 적었다. 계획을 고쳐 편차 없애기 — 가장 비싼 오류. 재기준선이라 부르지만 편차의 은폐. "17개월 계획이 20개월 걸렸는데 마지막 기준선 대비 편차 0" 보고서가 이렇게 만들어진다. 재기준선의 조건(3.4절 — CCB + 이사회, 최대 2회, 편차 이력 보존). 생각해 볼 질문 5를 던진다 — 재기준선 0회 프로젝트와 3회 프로젝트 중 어느 쪽이 잘 관리된 것인가, 답하려면 무엇을 더 알아야 하는가. 예상 소요 5분.}
\end{frame}

\begin{frame}[shrink=10]{EVM의 세 숫자와 그 앞의 하나 — PMB. ANSI/PMI 19-006-2019가 정의하는 것}
\begin{tbl}[\footnotesize]{L{0.2\textwidth}L{0.36\textwidth}L{0.36\textwidth}}
\textbf{표준 용어} & \textbf{정의} & \textbf{온담 P1} \\ \midrule
Control account & 범위·일정·원가가 통합되어 성과를 측정하는 단위 & WBS L2 다섯 산출물 + 원가 항목(SI·상용SW·인프라·단말·전환·시험·감리·PMO) \\
Work / Planning package & 측정 단위 / 아직 상세화되지 않은 원거리 예산 묶음 & 작업패키지 43개 / R4·R5 구간(rolling wave) \\
\textbf{PMB · BAC} & 시간대별로 배분된 예산의 합 — 관리예비비 제외 & \textbf{월별 PV 표 17개월, 합 145.80} \\
Undistributed budget & 승인되었으나 통제 계정에 미배분 & 승인 변경 CR-06·08·09의 1.005 — BL-02 편입 전까지 \\
Management reserve & PMB 밖, 미식별 리스크용, PM 권한 밖 & 관리예비비 12.0(정미란) \\
(Contingency) & 표준은 위치를 조직 관행에 맡긴다 & PMB 밖 · 집행 한도 안 · \textbf{집행 시 PV 없는 AC} \\
\end{tbl}
\keymsg{PV = 하기로 한 일의 예산 / EV = 한 일의 예산(계획 당시 값) / AC = 쓴 돈 — 셋이 같은 단위(돈)여서 일정과 원가가 한 표에 놓인다}
\note{§2.1 개념. "60\%의 기간이 지났는데 40\%의 가치가 인도되었다"는 문장은 시간과 가치가 같은 단위일 때만 성립한다 — EVM의 유일하고 결정적인 장점. 세 숫자 앞에 PMB가 있어야 한다 — Day2 4.3절이 월별 PV 표를 만든 이유. PMB 없으면 PV 없고, PV 없으면 SV·SPI 없다. 2019년 ANSI 표준은 2011년 Practice Standard 2판을 대체하고 적응형·증분 인도 개념을 통합한 현행 정본이다. 실무 오류 — BAC를 승인 예산 168.0으로 두어 CPI·TCPI가 좋아 보이는 것 / PMB 없이 계약 총액을 기간으로 균등 배분해 PV로 쓰는 것(4월 조달 봉우리 같은 구조가 사라져 편차 원인을 열별로 분리할 수 없다) / 관리예비비를 PMB에 넣는 것 / 컨틴전시 위치를 정하지 않아 집행이 어느 열의 AC로 나타나는지 매번 다른 것. 예상 소요 5분.}
\end{frame}

\begin{frame}{BAC는 145.8인가 156.0인가 — 온담이 실제로 부딪힌 질문}
\fig[0.56]{fig_2_1_pmb_bac_layers}
\figcap{그림 2-1. 컨틴전시 10.2는 PV 곡선 위 어디에도 배분되어 있지 않다. 온담은 BAC = 145.80(PMB), 예외 판정 기준 = 156.00(집행 한도)으로 분리했고 컨틴전시 집행은 PV 없는 AC로 기록한다 — CR-10이 헌장 문언으로 확정}
\keymsg{배분되지 않은 돈을 BAC에 넣으면 EV는 결코 BAC에 도달하지 못하고, SPI는 종료 시점에도 1이 되지 않으며, TCPI는 항상 실제보다 낮게 나온다}
\note{§2.1 "BAC는 145.8인가 156.0인가" — 워크북 ⑫ 해설 첫 항목. Day2 4.2절이 PMBOK 용어로 "원가 기준선 = 계획 원가 + 컨틴전시 = 156.0"이라 한 것은 PM이 책임지는 총액이라는 뜻이고 맞다. 그러나 EVM의 BAC는 시간대별로 배분된 PMB의 합이어야 하고 Day2 PV 표의 합은 145.80이다. 표준 용어로는 PMB(145.8)와 "관리예비비를 제외한 프로젝트 예산"(156.0)의 구분. 컨틴전시의 위치는 조직이 정한다 — PMB 안에 넣어 통제 계정별로 배분하는 조직도 있다. 온담은 밖에 두고 집행 시 PV 없는 AC로 기록하기로 했으므로 컨틴전시 집행 1.41 중 AC 발생분 1.19가 SI 열의 CV $-$1.19로 나타난다 — 원가 초과가 아니라 규칙의 결과이고, 그래서 SI 열의 원가 성과는 CPI가 아니라 "컨틴전시 소진율"로 읽는다(2교시 착시 2). CR-10이 이것을 헌장 §12 문언("초과 예상 판정 기준 = EAC가 집행 한도 156.0을 초과할 전망")으로 확정한다 — FP 0의 변경이 성과보고의 판정 기준을 바꾼 사례(3교시). 예상 소요 5분.}
\end{frame}

\begin{frame}[shrink=14]{Rules of Credit — 지표는 측정이 아니라 합의의 산물이다}
\fig[0.58]{fig_2_2_rules_of_credit}
\figcap{그림 2-2. 같은 작업, 여섯 규칙, 여섯 EV 곡선. 0/100·마일스톤 가중·물리적 \%·LOE가 온담 채택 규칙, 50/50과 배분 노력은 미사용(라이선스의 50/50은 AC 인식 규칙이지 EV 규칙이 아니다). 아래 두 상자 = 규칙의 어긋남이 만드는 구조적 편차(1.96 · 4.70)}
\keymsg{SPI 0.943은 자연에 있는 값이 아니라 G1에서 합의한 규칙이 만든 값이다 — 규칙이 달랐으면 다른 값이다}
\note{§2.2. 그림은 위 여섯 곡선을 왼쪽부터 읽고, 아래 두 상자로 내려온다. 여섯 기법과 온담 적용 — 0/100(설계 스프린트 리뷰 인수 1.25/회, 리허설 회차, 감리·시험 회차 보고서) / 50/50(미사용) / 마일스톤 가중(단말 입고 검수 70 · 설치 확인 30, 통합·시험 스프린트 = 통합시험 TC 통과율 60 · UAT 종료 40) / 물리적 \%(구현 스프린트 = 인수 FP ÷ 620 × 1.96, 전환 매핑 = 확정 테이블 수 ÷ 1,842) / 배분 노력(미사용) / LOE(PMO 월 0.56, 유지보수, 클라우드 — EV = PV). 기법의 선택이 SPI를 가른다 — 느슨한 규칙(착수 시 100\%, 담당자 진척률)은 SPI가 1 근처에 붙어 있다가 종료 직전 무너지고, 엄격한 규칙(진행 중 0)은 구조적으로 1 아래에서 시작해 실제 편차를 더한 값으로 움직인다. 두 가지 결론 — 규칙은 계획 시점에 확정하고 실행 중 바꾸지 않는다(바꾸면 EV는 원하는 숫자가 된다) / 문서로 남긴다(원가 기준선 v1.0 부속 2 [추가 설정], BL-02에서 정식 첨부 — 워크북 §6 갱신 6번). 실무 오류 — SPI가 나쁠 때 "진행 중 50\% 인정"으로 완화하는 것, LOE 항목을 물리적 \%로 재려는 것(LOE는 SV = 0이므로 AC로만 읽는다), 규칙을 남기지 않아 감리가 "EV 산정 근거"를 물을 때 답이 없는 것. 예상 소요 7분.}
\end{frame}

\begin{frame}[shrink=13]{PV 규칙과 EV 규칙의 시차 — 1.96 · EV 규칙과 AC 인식 기준의 시차 — 4.70}
\begin{itemize}
\item Day2 PV 배분 규칙 "스프린트 \textbf{시작 월}에 계상" vs EV 규칙 "스프린트 리뷰에서 \textbf{인수}된 FP를 인수 월에" — 각각 합리적(PV는 자금 조달 시점, EV는 인수 시점)
\item 둘이 어긋나면 월말에 걸친 스프린트 하나만큼 구조적 SV 음수 — 2026-09-30에 S20(09-28 착수)의 PV \textbf{1.96}. 정상 진행에서도 EV 0 \ra{} SV $-$4.40 중 $-$1.96은 지연이 아니다
\item 워크북 ⑫ 항목 4는 이것을 "조정 1: 측정 시차 제거 — SV $-$2.44, SPI 0.968"로 \textbf{별도 행}에 적는다. 병기이지 대체가 아니다
\item 같은 구조가 라이선스에 — EV 규칙 "인도·설치 확인서 100\%" vs AC 인식 "인보이스 발생액", 03-13 계약 변경으로 지급 조건이 인도 50\% / 통합시험 통과 50\% \ra{} EV 100\%, AC 50\% = CV \textbf{+4.70}
\item 이것은 EV 규칙의 문제가 아니라 \textbf{AC 인식 기준이 실행 중 바뀌었는데 양식에 조정 행이 없었던} 문제다
\end{itemize}
\keymsg{조정값을 본값 대신 쓰면 기준선 규칙을 실행 중에 바꾸는 것이다 — 옆에 적고, 이름을 붙이고, 크기를 적는다}
\note{§2.2 "온담의 사례로 읽기". 4교시 손계산에서 학생이 가장 자주 틀리는 자리 — 조정 SV를 구해 놓고 본값을 지운다. 미발생 4.70 = MDM 2차·APM 6.0 × 50\% + SAP 애드온 3.4 × 50\%(인도는 각각 5월·7월), 2027-01 통합시험 통과 시 한꺼번에 발생 — 그래서 2027-01 자금 봉우리(단말·DR·시험 16.20)에 얹혀 R-17의 일부가 된다. 학생에게 묻기 — "이재현 PM이 '진행 중 스프린트 50\% 인정'을 요구하면 무엇을 근거로 거절하는가"(제2장 생각해 볼 질문 4): 근거는 G1에 확정된 Rules of Credit 문서이고, 거절하지 않으면 SPI는 1 근처로 올라가 아무것도 재지 않게 된다. 예상 소요 4분.}
\end{frame}

\begin{frame}{편차와 지수 — 읽기 규칙 셋, 그리고 열별 분해는 의무다}
\begin{itemize}
\item 온담 2026-09-30 누적: PV 77.83 / EV 73.43 / AC 71.13 \ra{} SV $-$4.40 · CV +2.30 · SPI 0.943 · CPI 1.032
\item \textbf{SPI를 일수로 읽지 않는다} — 0.943은 비율. SV ÷ 월평균 PV 환산은 PV에 봉우리가 있으면 틀린다(4월 17.69 vs 9월 6.09, 세 배). 일수는 ES(2.4절) 또는 CPM(2.8절)이 준다
\item \textbf{월별 지수와 누적 지수를 섞지 않는다} — 누적 SPI 0.592 \ra{} 0.796 \ra{} 0.898 \ra{} \textbf{0.784(4월)} \ra{} 0.871 \ra{} 0.881 \ra{} 0.914 \ra{} 0.913 \ra{} 0.943. 회복의 일부는 수렴 착시
\item \textbf{열별로 분해한다}:
\end{itemize}
\begin{tbl}[\scriptsize]{L{0.08\textwidth}rrrrrrr}
 & \textbf{SI} & \textbf{상용SW} & \textbf{인프라} & \textbf{전환} & \textbf{시험·감리} & \textbf{PMO} & \textbf{합} \\ \midrule
SV & $-$5.31 & 0 & 0 & +0.20 & +0.70 & 0 & $-$4.40 \\
CV & $-$1.19 & \textbf{+4.70} & $-$1.00 & +0.20 & $-$0.05 & $-$0.36 & +2.30 \\
\end{tbl}
\keymsg{합계는 "일정은 조금 늦고 원가는 남는다" — 열별은 "개발이 늦고 감리 선행이 가리며, 원가는 라이선스 이연이 인프라·PMO 초과를 덮는다"}
\note{§2.3 "편차와 지수". 4월 0.784는 상용SW $-$6.0(조달 이연)이 SV $-$7.97의 75\% — 균등 배분 곡선이었다면 4월 SPI는 0.9쯤이었을 것이고 "약간의 지연"으로 넘어가 5월 R-05와 합쳐져 "지난달부터 조금씩 늦었다"는 잘못된 서사가 되었을 것이다. 봉우리 덕분에 4월은 조달(IS-01), 5월은 설계(IS-02)라는 두 사건으로 분리되었고 각각 다른 대응(지급 조건 협상 / 대체 결정권자)을 갖는다 — Day2 §4.6 해설의 예고가 실현된 자리. 표의 SI CV $-$1.19는 컨틴전시 집행 4건(R-07 0.15 + R-05 0.42 + R-03 0.50 + CR-03 0.12)이고, 인프라 $-$1.00은 클라우드 실사용(월 0.4 \ra{} 0.5, IS-10), PMO $-$0.36은 인건비·도구. 예상 소요 5분.}
\end{frame}

\begin{frame}{1교시 핵심 정리}
\begin{enumerate}
\item 통제 = 측정(같은 규칙) \ra{} 분석(원인) \ra{} 예측(전제) \ra{} 조치(세 층 — 격자가 가른다) \ra{} 기준선 갱신(승인된 변경만). 실무적 정체는 보고서 양식이다
\item 하이브리드 통제는 스프린트 / 월 / 게이트 세 층이고, 연결 규칙(Rules of Credit · 접수 마감 · 캘린더 = CPM 날짜)이 없으면 아래층의 작은 결정이 위층의 큰 사건이 될 때까지 안 보인다
\item PMBOK 8판 두 Focus Area · 다섯 성과영역, PRINCE2 7 세 프로세스 · 예외에 의한 관리(예외 보고는 위임의 반대급부), ISO 21502 §7(이슈와 리스크 별도), 한국 SI 네 주기 · "감리 지적은 검수의 근거"
\item 비싼 오류 셋 — 90\% 증후군 · 기준선 없는 보고 · 재기준선 남용 — 의 공통점은 기록의 부재
\item BAC = PMB의 합 145.80, 예외 기준 = 집행 한도 156.00, 컨틴전시는 PMB 밖·집행 시 PV 없는 AC. EV는 합의의 산물(Rules of Credit) — 구조적 시차 1.96·4.70은 병기하되 대체하지 않는다. SPI는 일수가 아니고, 열별 분해는 의무
\end{enumerate}
\keymsg{2교시는 "이대로 가면 어디에 도착하는가" — EAC 4공식, Earned Schedule, 다섯 착시, 그리고 임계경로 $-$10}
\note{교재 제1장 핵심 정리 1\til{}5와 제2장 핵심 정리 1\til{}2, 3의 전반. 생각해 볼 질문(제1장 1 — 여러분 회사 월간 보고서의 "진척률"은 무엇의 비율인가, 보고서 안에 정의가 있는가 / 제1장 3 — 예외 보고를 권한의 반납으로 느끼는 PM과 위임의 반대급부로 느끼는 PM의 행동은 어떻게 다른가, 여러분 조직에서 예외 보고를 낸 PM은 어떤 대우를 받았는가) 중 하나를 쉬는 시간 전에 던진다. 예상 소요 3분.}
\end{frame}

% ===== 2교시 (90분) — EAC · Earned Schedule · 착시 · 흐름 지표 · 임계경로 회복 =====
\section{2교시 (90분) — EAC · Earned Schedule · 착시 · 흐름 지표 · 임계경로 회복}

% =====================================================================
% 2교시 — EAC · Earned Schedule · 착시 · 흐름 지표 · 임계경로 회복
% =====================================================================
\begin{frame}{2교시에서 배우는 것 — "이대로 가면 어디에 도착하는가"}
\begin{itemize}
\item \textbf{EAC 4공식}은 4개의 전제 — 공식 EAC는 근거가 있는 것 하나(EAC4 150.86), 검산은 EAC1$'$, 경고는 EAC3$'$
\item \textbf{TCPI}는 분모를 적어야 숫자다 — BAC / 조정 AC / 집행 한도 / EAC4
\item \textbf{Earned Schedule} — SPI는 종료로 수렴한다. 시간으로 재면 SV(t) $-$0.72개월, IEAC(t) 18.48개월
\item \textbf{착시 5개}에 이름·크기·조정값을 붙여 병기한다 — 대체하지 않는다
\item \textbf{흐름 지표}는 EVM의 2주 앞 신호, \textbf{CPM}은 청진기 — A10 $-$10이 EX-01을 만든다
\end{itemize}
\keymsg{온도계(EVM) \ra{} 보정표(ES) \ra{} 청진기(CPM) — 세 도구가 같은 환자를 다르게 말할 때 어느 것을 믿는가}
\note{교재 제2장 2.3 후반\til{}2.9, 워크북 §2(⑫ 성과보고) 항목 4\til{}9. 1교시가 "지금 어디 있는가"(편차·지수)였다면 2교시는 "이대로 가면 어디에 도착하는가"(예측)와 "왜 지표가 실제와 다르게 보이는가"(착시)이고, 마지막에 "그래서 누가 무엇을 결정하는가"(EX-01)로 끝난다. 모든 숫자는 워크북 완성 예시본 — 4교시 손계산에서 학생이 같은 값을 다시 만들게 된다. 도구 비유(온도계·보정표·청진기)를 첫 프레임에서 세워 두고 교시 내내 반복한다. 예상 소요 3분.}
\end{frame}

\begin{frame}{EAC 4공식 — 어느 공식이 맞는가가 아니라 어느 전제가 맞는가}
\begin{tbl}[\scriptsize]{L{0.30\textwidth}L{0.30\textwidth}rrL{0.18\textwidth}}
\textbf{공식} & \textbf{전제} & \textbf{EAC} & \textbf{VAC} & \textbf{판정} \\ \midrule
EAC1 = BAC ÷ CPI & 누적 원가 효율 유지 & 141.23 & +4.57 & 기각 — 착시 포함 \\
EAC2 = AC + (BAC $-$ EV) & 편차는 일회성 & 143.50 & +2.30 & 기각 — 4.70 미반영 \\
EAC3 = AC + (BAC $-$ EV) ÷ (CPI×SPI) & 원가·일정 효율 동시 & 145.43 & +0.37 & 참고 \\
\textbf{EAC4 = AC + 상향식 ETC} & 항목별 잔여 재추정 & \textbf{150.86} & \textbf{$-$5.06} & \textbf{공식} \\
EAC1$'$ = BAC ÷ 조정 CPI 0.968 & 발생주의 조정 후 유지 & 150.57 & $-$4.77 & EAC4 검산 \\
EAC3$'$ = 조정 AC + 잔여 ÷ (0.968×0.943) & 조정 CPI × SPI & \textbf{155.04} & $-$9.24 & \textbf{경고 대역} \\
\end{tbl}
\keymsg{PMB 145.8 아래의 EAC1은 고정가 SI 68.2가 줄지 않는 한 구조적으로 불가능하다 — 공식이 아니라 전제를 기각한다}
\note{§2.3 "EAC 4공식과 각각의 가정", 워크북 ⑫ 항목 5 표. 네 공식을 다 적는 이유는 보고서 안에서 전제를 논증하기 위해서다 — EAC 하나만 적는 성과보고는 대개 가장 낙관적인 EAC1을 고르고 스폰서는 원가가 남는다고 믿는다. EAC1 141.23은 CPI 1.032(라이선스 이연 착시 포함)의 외삽이고, EAC2 143.50은 미발생 4.70과 컨틴전시 잔여 집행이 잔여에 없다. EAC3$'$ 155.04는 일정 회복 실패 시의 도착지로 집행 한도까지 0.96 — 스폰서가 기억해야 할 두 번째 숫자(첫 번째는 A10 $-$10). 4교시에서 학생은 이 표를 손으로 다시 만든다. 예상 소요 6분.}
\end{frame}

\begin{frame}{EAC4의 상향식 ETC 79.72 — 행마다 근거가 있어야 "공식"이 된다}
\begin{itemize}
\item SI 잔여 계약 38.51 + 변경 대가 1.02 + 대기 잔여 0.18 — 근거: 계약 잔액·변경 로그
\item 상용SW 5.90 = 잔여 PV 1.20 + \textbf{미발생 4.70} — 근거: 지급 조건 50/50, 2027-01 통합시험 통과 시 발생
\item 인프라 7.92 = 7.40 × 1.07 — 근거: 클라우드 실사용 추세(IS-10)
\item 단말 8.90 · 전환 5.10 · 시험·감리 7.40 · PMO 4.80 — 근거: 견적·잔여 PV
\item 검산: 조정 CPI의 EAC1$'$ 150.57과 EAC4 150.86이 \textbf{0.3 안}에서 만난다 — 두 방법이 크게 다르면 어느 전제가 틀린 것
\end{itemize}
\keymsg{수행사 견적은 계약 잔액이지 발주사의 잔여 원가가 아니다 — 상향식 ETC를 견적으로 대체하지 않는다}
\note{§2.3 같은 절. 상향식 ETC의 각 행에 근거 열이 있다는 것이 EAC4가 "공식"이 되는 조건이다. 상용SW 행의 4.70이 여기 들어가므로 EAC4는 착시 1을 이미 걷어 낸 값이고, 그래서 착시를 걷어 낸 EAC1$'$과 만난다 — 만난다는 것 자체가 검산이다. 인프라 1.07은 월 0.4 \ra{} 0.5 실사용 추세를 잔여 7개월에 적용한 것. 실무 오류 — 수행사 견적으로 ETC를 대체하는 것(발주사 관점 잔여 원가에는 단말·인프라·PMO가 있다), 변경 대가 1.02를 ETC에 넣지 않는 것(CR-08\til{}11 미승인이라도 전망에는 넣고 "미반영 변경"으로 표시). 예상 소요 4분.}
\end{frame}

\begin{frame}{TCPI 4종과 VAC — 분모를 적지 않은 TCPI는 숫자가 아니다}
\begin{tbl}[\scriptsize]{L{0.30\textwidth}L{0.28\textwidth}rL{0.32\textwidth}}
\textbf{TCPI} & \textbf{산식} & \textbf{값} & \textbf{읽기} \\ \midrule
TCPI(BAC) & 72.37 ÷ 74.67 & 0.969 & 본값 — 착시 포함, 의미 없음 \\
TCPI(BAC, 조정 AC 75.83) & 72.37 ÷ 69.97 & \textbf{1.034} & 잔여 3.4\% 절감 필요 \ra{} PMB 초과 확정 \\
TCPI(집행 한도 156.00) & 72.37 ÷ 84.87 & 0.853 & 집행 한도 안은 가능 — 컨틴전시가 완충 \\
TCPI(EAC4 150.86) & 72.37 ÷ 79.73 & 0.908 & 공식 EAC 달성에 필요한 효율 \\
\end{tbl}
\begin{itemize}
\item VAC = BAC $-$ EAC4 = $-$5.06 · 집행 한도 여유 5.14 \ra{} \textbf{원가 예외 보고 불요}
\item 단 EAC3$'$ 155.04(여유 0.96) — 10-31 데이터에서 EAC4 > 153.0이면 예외 보고 준비를 \textbf{예고}(경고 A)
\item 승인 예산 168.0과 비교한 여유 17.14는 관리예비비 12.0을 포함한 남의 돈
\end{itemize}
\keymsg{"TCPI > 1.1이면 회복 불가"를 인용할 때 분모가 BAC인지 집행 한도인지 먼저 — 온담은 BAC로는 초과 확정, 한도로는 여유}
\note{§2.3 "TCPI와 VAC", 워크북 ⑫ 항목 5 하단 표. 1교시의 BAC 논쟁(145.8 vs 156.0)이 여기서 실제 결과로 갈린다 — 같은 잔여 작업 72.37을 어느 돈으로 끝내는가에 따라 1.034와 0.853. 원가 RAG가 R이 아니라 A인 이유가 이 두 줄이다(집행 한도 기준 여유 5.14, 그러나 EAC3$'$ 경고). 예고 문장("다음 보고에서 153.0을 넘으면 예외 보고를 준비한다")은 PRINCE2의 경고 등급 운용 — 예외 보고를 갑자기 내지 않기 위한 장치다. 실무 오류: VAC를 승인 예산과 비교하는 것, CV 양수를 절감으로 읽는 것 — 설명할 수 없는 양수 CV는 대개 AC 인식의 지연. 예상 소요 5분.}
\end{frame}

\begin{frame}{EAC 사다리 — 스폰서가 기억할 숫자는 둘이다}
\fig[0.58]{fig_2_3_eac_ladder}
\figcap{그림 2-3. 온담 2026-09-30의 EAC 사다리. 네 공식은 네 전제이고, 항목별 근거를 가진 EAC4(150.86)만 공식이 될 수 있으며 조정 CPI의 EAC1$'$(150.57)이 검산한다. EAC3$'$(155.04)은 회복 실패 시 도착지 — 집행 한도까지 0.96}
\keymsg{승인 168.0 / 집행 한도 156.0 / PMB 145.8 세 층 사이에 EAC 여섯 값을 놓으면 "원가는 A"가 한눈에 읽힌다}
\note{§2.3 그림 2-3. 사다리를 아래에서 위로 읽는다 — EAC1 141.23(PMB 아래, 불가능) \ra{} EAC2 \ra{} EAC3 \ra{} PMB 145.8 \ra{} EAC1$'$·EAC4(150.6\til{}150.9, 공식) \ra{} EAC3$'$ 155.04 \ra{} 집행 한도 156.0 \ra{} 승인 168.0. 오른쪽 TCPI 4값이 각 분모에 대응한다. 학생에게 묻기 — "스폰서 정미란에게 한 문장으로 원가를 보고하라": 정답은 "공식 EAC 150.86으로 PMB를 5.06 넘지만 집행 한도 안(여유 5.14)이며, 일정 회복이 실패하면 155.04까지 간다"이다. 예상 소요 3분.}
\end{frame}

\begin{frame}{SPI의 결함 — 6개월 늦게 끝난 프로젝트의 마지막 SPI도 1.0이다}
\begin{itemize}
\item EV는 결국 BAC에 도달하고 PV도 계획 종료 후 BAC에 머문다 \ra{} SPI = EV ÷ PV는 \textbf{종료 시점에 반드시 1.0}
\item 수렴은 종료 순간이 아니라 후반부 내내 진행 — 마지막 1/3에서 SPI는 실제 지연과 무관하게 \textbf{좋아지는 것처럼} 보인다
\item SV(돈 단위)도 같은 이유로 0으로 수렴 — Lipke(2003) 「Schedule is Different」
\item 해법: 일정 편차를 돈이 아니라 \textbf{시간}으로 잰다 — Earned Schedule
\item ES = "현재 누적 EV를 PV 곡선이 계획상 \textbf{언제} 달성하기로 했는가"의 시점
\end{itemize}
\keymsg{후반부 프로젝트에서 SPI를 보고 안심하는 조직이 겪는 일 — 온담은 진척 50\%에서 이미 벌어지기 시작했다}
\note{§2.4 "결함". 칠판에 EV·PV 두 곡선을 그려 종료 직전 두 곡선이 같은 BAC 수평선에 붙는 모습을 보이면 결함이 한 번에 이해된다. 학생 대부분이 SPI를 "일정 효율"로만 알고 있고, 종료 수렴을 처음 듣는다 — 여기서 "여러분 조직의 마지막 월간 보고서 SPI가 얼마였는가"를 묻는다. 대개 1.0 근처이고 그 프로젝트가 정시에 끝났는지 물으면 대답이 갈린다. 예상 소요 4분.}
\end{frame}

\begin{frame}{Earned Schedule — ES · SV(t) · SPI(t) · IEAC(t), 온담 2026-09-30}
\begin{tbl}[\scriptsize]{L{0.22\textwidth}L{0.40\textwidth}L{0.30\textwidth}}
\textbf{지표} & \textbf{산식 · 온담 값} & \textbf{읽기} \\ \midrule
ES & C + (EV $-$ PV$_C$) ÷ (PV$_{C+1}$ $-$ PV$_C$) = 8 + 1.69 ÷ 6.09 = \textbf{8.28개월} & 8월 71.74 < EV 73.43 < 9월 77.83 \ra{} C = 8 \\
SV(t) & ES $-$ AT = 8.28 $-$ 9 = \textbf{$-$0.72개월} & ≈ $-$15영업일(평균) \\
SPI(t) & ES ÷ AT = \textbf{0.920} & SPI 0.943과 0.023 차이 \\
IEAC(t) & PD ÷ SPI(t) = 17 ÷ 0.920 = \textbf{18.48개월} & +1.48개월 ≈ +31영업일 \\
추세 종료 & \textbf{2027-07-12} vs 허용범위 06-18(+15영업일) & \textbf{초과(추세)} \\
\end{tbl}
\begin{itemize}
\item IEAC(t)는 추세 외삽이지 계획이 아니다 — 회복 계획(to-be MG3 02-19)과 나란히 놓고 어느 것이 계획인지 적는다
\item ES는 프로젝트 전체의 \textbf{평균} — 지연은 임계경로 A10에 $-$10으로 몰려 있다
\end{itemize}
\keymsg{ES는 온도계의 눈금을 바로잡는 보정표이지 청진기가 아니다 — 청진기는 CPM}
\note{§2.4 "Earned Schedule", 워크북 ⑫ 항목 6. 보간 계산을 학생과 함께 한 번 한다 — PV 누적 표(Day2 §4.5)에서 EV 73.43을 넘지 않는 마지막 달(8월 71.74)을 찾고 9월 증분 6.09로 나눈다. AT = 9(1월 착수, 9월 말). PD 17은 기준선 계획 기간(2027-05 종료). +31영업일은 1.48개월 × 약 21영업일. 실무 오류 — SPI만 보고 "회복 중"이라고 읽는 것, ES 없이 SV를 일수로 환산하는 것(1교시 봉우리 문제), IEAC(t)를 종료일로 보고하는 것, ES를 계산하고 CPM을 보지 않는 것. 워크북 항목 6의 추세 도착일 표기가 두 곳(07-12 / 07-14)으로 다르다 — 성과보고 요약표의 07-12를 정본으로 쓰고 4교시에서 정정한다. 예상 소요 7분.}
\end{frame}

\begin{frame}{SPI vs SPI(t) 월별 — 7월부터 벌어진다}
\begin{tbl}[\scriptsize]{lrrrL{0.46\textwidth}}
\textbf{월} & \textbf{SPI} & \textbf{SPI(t)} & \textbf{차이} & \textbf{판독} \\ \midrule
2026-04 & 0.784 & 0.887 & +0.103 & 조달 이연 — 돈으로 $-$8.0, 시간으로 $-$0.45개월 \\
2026-05 & 0.871 & 0.868 & $-$0.003 & 라이선스 인도로 SPI 회복, SPI(t)는 R-05 반영 \\
2026-07 & 0.914 & 0.904 & $-$0.010 & \textbf{벌어지기 시작} \\
2026-08 & 0.913 & 0.914 & +0.001 & \\
\textbf{2026-09} & \textbf{0.943} & \textbf{0.920} & \textbf{$-$0.023} & SPI는 회복처럼, SPI(t)는 정체 — 수렴 착시 \\
\end{tbl}
\begin{itemize}
\item 9월 SPI는 8월보다 0.030 상승, SPI(t)는 0.006 상승 — 지수로는 회복, 시간으로는 정체
\item 4월의 역방향(+0.103)은 PV 봉우리 — 돈으로 큰 조달 이연이 시간으로는 0.45개월
\end{itemize}
\keymsg{진척 50\%에서 이미 0.023 벌어졌다 — 후반부에는 더 벌어진다. 일정 보고는 SPI(t)·ES·CPM TF로}
\note{§2.4 월별 표(3월·6월 행은 화면에서 생략, 교재 표에 있다). 두 방향의 차이를 다 보인다 — 봄에는 SPI(t)가 관대(ES가 PV 완만 구간에 있어서), 여름부터는 SPI가 관대(수렴). 즉 어느 쪽이 항상 보수적인 것이 아니라 PV 곡선의 모양과 진척 단계에 따라 다르고, 그래서 둘을 병기한다. 5월 SPI 회복(0.784 \ra{} 0.871)이 라이선스 인도(EV +6.0)로 생긴 것이지 개발 회복이 아니라는 점은 열별 분해(1교시)와 연결된다. 예상 소요 4분.}
\end{frame}

\begin{frame}{Earned Schedule 그림 — EV 73.43이 PV 곡선에서 도달한 시점}
\fig[0.58]{fig_2_4_earned_schedule}
\figcap{그림 2-4. 위 — EV 73.43이 PV 곡선에서 도달한 시점(ES 8.28개월)과 경과 기간(AT 9)의 간격이 SV(t) $-$0.72개월. 아래 — SPI는 종료로 갈수록 1로 수렴하고 SPI(t)는 그렇지 않다. 온담에서 둘은 7월부터 벌어지기 시작했다. Lipke(2003)를 워크북 ⑫ 항목 6의 값으로 재구성}
\keymsg{위 그림의 수평 간격이 일정 편차다 — 수직 간격(SV $-$4.40)은 돈이지 시간이 아니다}
\note{§2.4 그림 2-4. 위 패널에서 EV 수평선이 PV 곡선과 만나는 점을 x축으로 내려 8.28을 읽는 동작을 레이저로 보인다 — SV(수직)와 SV(t)(수평)의 차이가 한 그림에 있다. 아래 패널의 두 선이 7월부터 갈라지는 자리를 짚고, 앞으로 갈수록 위쪽 선(SPI)만 1로 올라갈 것임을 손으로 그려 준다. 예상 소요 3분.}
\end{frame}

\begin{frame}{CPI 안정성 — Christensen 규칙이 말하는 것과 말하지 않는 것}
\begin{itemize}
\item 연구: 미 국방 획득 사업 대규모 표본에서 누적 CPI가 \textbf{진척 20\% 이후 ±10\% 안에서 안정} — "CPI는 대개 좋아지지 않는다"
\item 말하지 않는 것 ① \textbf{표본이 다르다} — 원가 정산 계약 vs 온담 고정가 SI(그 열의 CPI는 계약이 만든 1)
\item ② \textbf{안정된 값이 무엇인지가 먼저} — 온담 누적 CPI 0.978 \ra{} 0.994 \ra{} \textbf{1.073(5월)} \ra{} 1.032은 03-13 지급 조건 변경의 흔적. 조정 0.968이 안정값
\item ③ \textbf{진척 50\%는 20\%가 아니다} — 규칙이 성립한다면 0.968은 크게 회복되지 않는다 \ra{} PMB 초과(TCPI 1.034)는 확정
\item 착시 포함 CPI에 규칙을 적용하면 EAC1 141.23 — PMB보다 낮은, 고정가 구조에서 불가능한 숫자
\end{itemize}
\keymsg{Christensen 규칙은 "CPI가 좋아질 것"이라는 희망을 꺾는 데 쓰는 것이지, 착시 포함 CPI를 외삽하는 데 쓰는 것이 아니다}
\note{§2.5. Christensen \& Payne(1992, Journal of Parametrics)·Christensen \& Heise(1993, NCMJ) — 서지와 후속 논쟁은 6.3절(마무리 교시). 규칙의 실무 함의는 명확하다("20\%에서 CPI 0.9면 곧 회복된다는 말을 믿지 말라")지만 표본·계약 구조·착시 조정 세 조건을 확인하고 쓴다. 온담 5월의 1.073 점프를 학생에게 "원가 효율이 좋아진 것인가"로 물으면 대부분 그렇다고 답한다 — 착시 1의 예고편. 예상 소요 4분.}
\end{frame}

\begin{frame}{흐름 지표 — EVM이 못 보는 것: throughput · WIP · work item age}
\fig[0.50]{../그림/PM_Day3/fig_2_7_flow_metrics.pdf}
\figcap{그림 2-7. S7\til{}S18 throughput(FP/스프린트) 620·620·600·480·470·560·560·600·620·620·640·630, WIP 최고 41(S10), work item age 최장 12영업일}
\keymsg{S10\til{}S11의 470\til{}480은 SPI에 두 달 뒤 나타났다 — 흐름 지표는 EVM의 선행 지표다}
\note{교재 2.7 "흐름 지표", 워크북 §2.5 항목 9. Vacanti의 네 지표를 온담 스프린트 데이터로 읽는다. throughput 하락(S10 480·S11 470)은 R-05 설계 지연이 구현 스프린트로 번진 시점이고, 같은 사실이 월별 SPI에는 5\til{}6월(0.871·0.881)에 완만하게만 나타난다 — 스프린트 단위 지표가 월 단위 EVM보다 먼저 움직인다. WIP 41(S10)은 설계 승인 대기가 쌓인 것이고, work item age 12일은 "완료된 것만 재는" cycle time에서는 보이지 않는다. "완료 기준 cycle time만 개선되는 착시"와 SPI 수렴 착시는 같은 병리라는 점을 학생에게 묻는다(둘 다 끝난 것만 세기 때문이다). CFD는 그림의 하단 — 진행 중 띠의 두께가 WIP다. 예상 소요 5분.}
\end{frame}

\begin{frame}{DORA 4지표의 위치 — 배포 파이프라인의 건강, 프로젝트 성과와는 다른 층}
\begin{itemize}
\item 온담 P1 [추가 설정] — 배포 빈도 주 2회 · 변경 리드타임 3.1일 · 변경 실패율 12\% · 복구 4시간
\item DORA는 \textbf{릴리스 파이프라인}을 재고, EVM은 \textbf{기준선 대비 편차}를 잰다 — 둘은 대체 관계가 아니다
\item 실패율 12\%는 8월 시험 정지(TC 210 \ra{} 95)와 같은 시기 — 시험이 줄면 배포 실패가 늘어난다
\item 하이브리드의 세 층: 스프린트(흐름·DORA) / 월(EVM·ES) / 게이트(허용범위·예외)
\item 한 층의 지표를 다른 층에 올리지 않는다 — 스토리포인트는 EV가 아니고, 배포 빈도는 진척이 아니다
\end{itemize}
\keymsg{세 층은 각자 다른 질문에 답한다 — "잘 흐르는가 / 계획대로인가 / 허용범위 안인가"}
\note{교재 2.7 후반. DORA 값은 회사 정본에 없어 [추가 설정]으로 워크북 §2.5 항목 9에 두었다. 강조할 것은 층의 분리다 — 현장에서 가장 흔한 혼동은 velocity(스토리포인트/스프린트)를 진척률로 보고하는 것이다. 온담은 FP 인수 기준(Rules of Credit)으로 EV를 재므로 velocity와 EV가 분리되어 있고, 그것이 이 프로젝트 성과 측정의 강점이다. 변경 실패율 12\%가 8월 시험 인력 이동과 겹친다는 관찰은 5교시 결함 곡선으로 이어진다. 예상 소요 3분.}
\end{frame}

\begin{frame}{임계경로 — A10 $-$10영업일이 컷오버 창을 밀어낸다}
\fig[0.50]{../그림/PM_Day3/fig_2_8_cpm_asis_tobe.pdf}
\figcap{그림 2-8. as-is: A10 착수 09-28(기준선 180 \ra{} 실제 190) \ra{} MG3 2027-03-05 \ra{} 1차 컷오버 창(02-26) 이탈 \ra{} 2차 창 03-12\til{}14 \ra{} G4 06-11(+10). 합류 버퍼 15 전부 소진}
\keymsg{일정 편차는 돈이 아니라 달력에 나타난다 — SV $-$4.40은 "컷오버 창을 놓친다"는 문장으로 번역되어야 결정이 된다}
\note{교재 2.8 "임계경로 소진", 워크북 §2.5 항목 10. Day2 CPM 표의 A10(구현 R3, 기준선 착수 작업일 180)이 R2 종료 지연(+10)으로 190에 착수했고 임계경로라 그대로 MG3(G3 게이트)를 10영업일 민다. 1차 컷오버 창 02-26\til{}28은 주말 고정이므로 다음 창 03-12\til{}14로 넘어가고, 이것이 R-14(두 번 컷오버)·R-17(2027 자금 집중)을 만든 원인이다. Day2가 남겨 둔 합류 버퍼 15영업일은 여기서 전부 쓰인다 — "버퍼는 누구의 것인가"(Day2 §3 여유 소유권)가 실전이 되는 순간. 학생에게 as-is 그림만 보이고 "무엇을 결정해야 하는가"를 먼저 묻고 다음 프레임의 3대안으로 간다. 예상 소요 5분.}
\end{frame}

\begin{frame}{회복 계획 3대안과 예외 보고 EX-01 — 10-02 제출 · 10-16 결정 시한}
\begin{tbl}[\scriptsize]{L{0.6cm}L{4.6cm}L{2.6cm}L{3.4cm}}
 & 내용 & 효과 & 대가·조건 \\ \midrule
1 & A14 $-$5(SS 겹침) + A16 $-$5(UAT 병행) + R2 이월 420 \ra{} R3 120·A12 300 + CR-08 A12 편입 & MG3 02-19 회복, 여유 0 & A12 TF 10 \ra{} 0 — 제2 임계경로 생성. \textbf{권고} \\
2 & 2차 컷오버 창 03-12\til{}14 수용 & 압축 없음 & 관리예비비 +1.6(두 번 컷오버 운영비) \\
3 & 옴니 "매장 간 이동 요청" 약 260FP \ra{} R5 이연 & A10 부하 감소 & 범위 축소 — 헌장 §5 문언 · 조건부(11-13 S22 속도 620 미만 시 발동) \\
\end{tbl}
\keymsg{예외 보고는 사고 보고가 아니다 — "허용범위 안에서는 보고하지 않을 권리"의 반대급부이며, 대안과 권고를 실어야 문서가 된다}
\note{교재 2.8 후반과 1.3(예외에 의한 관리), 워크북 §2.5 항목 10\til{}11. EX-01의 구조 — 편차(임계 $-$10, IEAC(t) 18.48 vs 허용 종료 06-18) · 원인(R-05 실현·R2 이월 420FP) · 대안 3 · 권고(대안 1 + 대안 3 조건부) · 결정 시한(10-16, 스폰서). 대안 1의 대가가 "A12의 여유 10을 0으로 만든다"는 점을 학생이 스스로 찾게 한다 — 압축은 항상 다른 경로의 여유를 먹는다. 대안 2는 압축이 아니라 수용이고 대가가 관리예비비(스폰서 결재)라 PM 권한 밖이다. 대안 3은 범위 축소이므로 헌장 문언 변경 = CCB 필수(3교시 규칙 ④). 10-16 결정 시한은 S20 착수(10-12) 뒤라는 점 — 결정이 늦으면 대안 1의 A14 SS 겹침이 불가능해진다. 예상 소요 6분.}
\end{frame}

\begin{frame}{2교시 핵심 정리}
\begin{enumerate}
\item EAC 4공식은 네 개의 \textbf{전제}다 — 스폰서가 기억할 숫자는 둘(EAC4 150.86 · EAC3′ 155.04, 경고 대역)
\item Earned Schedule — ES 8.28 · SPI(t) 0.920 · IEAC(t) 18.48개월 \ra{} 2027-07-12 > 허용 06-18. SPI 0.943은 수렴 착시의 시작
\item 착시 셋을 걷어내면 CPI 1.032 \ra{} 0.968 — "원가가 남는다"는 라이선스 이연이 만든 문장
\item 흐름 지표는 EVM의 선행 지표 — S10\til{}S11 throughput 470이 두 달 뒤 SPI에 나타났다
\item 편차는 결정 요청으로 끝나야 한다 — EX-01, 대안 3, 권고 1, 시한 10-16
\end{enumerate}
\keymsg{3교시는 "그 결정을 문서로 어떻게 통과시키는가" — 변경요청서·6축·CCB·기준선 갱신}
\note{교재 제2장 핵심 정리 3\til{}8. 생각해 볼 질문(제2장 2 — "여러분 회사의 마지막 성과보고에서 SPI(t)를 계산하면 어떻게 되는가")을 던지고 답은 받지 않는다. 오후 4교시에서 이 숫자들을 손으로 다시 만든다는 것을 예고. 예상 소요 4분.}
\end{frame}

% ===== 3교시 (90분) — 변경 관리: 변경은 사건이 아니라 절차다 =====
\section{3교시 (90분) — 변경 관리: 변경은 사건이 아니라 절차다}

% =====================================================================
% 3교시 — 변경 관리 (교재 제3장, 워크북 §1)
% =====================================================================
\begin{frame}{3교시 — 이 교시에서 배우는 것}
\begin{itemize}
\item 변경의 원천 4종과 변경요청서의 3단 구조(요청서 \ra{} 영향 분석 \ra{} 의결서)
\item \textbf{6축 영향 분석} — FP 하나로는 변경의 크기를 잴 수 없다
\item 처리 경로 6규칙 — 핵심은 ② 헌장 문언 변경은 허용범위 내라도 CCB
\item 기준선 갱신 BL-02 vs 재기준선 — "반영한 뒤 편차가 줄어드는가"
\item 변경 대가 55만 원/FP · 감액도 변경 · 범위 크립의 통제 도구 셋
\end{itemize}
\keymsg{온담 변경 4건(CR-08\til{}11)이 같은 문을 통과해 BL-02가 되는 과정을 따라간다 — 4교시 실습 ①의 이론}
\note{교재 제3장 전체, 워크북 §1. Day2 2.4절이 범위 팽창의 세 형태와 처방(FP 기준선 확인서·변경 대가 경로)을 세웠다면, 이 교시는 그 경로를 실제로 통과시키는 절차를 다룬다. 이 교시의 숫자는 모두 워크북 §1.5 완성 예시본에서 가져왔고 4교시 실습 ①에서 수강생이 같은 숫자를 손으로 다시 쓴다. 강의 순서는 원천과 구조(3.1) \ra{} 6축(3.2) \ra{} 경로(3.3) \ra{} BL-02와 재기준선(3.4) \ra{} 대가와 크립(3.5) \ra{} 온담 4건과 CCB 제10차(3.6) \ra{} 핵심 정리다. 예상 소요 3분.}
\end{frame}

\begin{frame}{변경은 사건이 아니라 절차다 — 원천 4종}
\begin{tbl}[\scriptsize]{L{2.1cm}L{4.6cm}L{4.2cm}}
원천 & 뜻 & 온담 2026-09-30 \\
\midrule
요구의 변경 & 새 것을 원하거나 있던 것을 바꾼다 & CR-09 22:00 마감(서정필) \\
오류의 정정 & 기준선이 틀렸음을 발견 & CR-08 임시 연동 누락(WBS 1.1.7 "FP 0") \\
외부 사건 & 규제·벤더·타 프로젝트가 바뀜 & CR-01 배달3사 규격 v2 · CR-11 보존 요건 \\
거버넌스 변경 & 규칙·권한·정의 자체가 바뀜 & CR-10 원가 허용범위 정의(FP 0) \\
\end{tbl}
\keymsg{"요구가 바뀌었다"는 사건이다 — 그것이 기준선을 바꾸는가는 절차가 결정한다}
\note{교재 3.1. 원천이 다르면 영향 분석의 무게가 다르다 — 요구 변경은 대안 3개가 핵심이고, 오류 정정은 "왜 몰랐는가"가 등록부(R-06 OD-POS 지식)로 이어지며, 외부 사건은 귀책 분해(5교시)와 붙고, 거버넌스 변경은 FP가 0이어도 CCB다. 절차가 없으면 사건이 곧 기준선 변경이 되고 그것이 범위 크립의 정의다. 2026-09-30 시점의 변경 4건이 네 원천을 하나씩 대표한다는 점을 강조한다 — 시나리오가 그렇게 설계되었다. 실무 오류 — 요청자 칸에 "현업"·"PM"이 적히는 것(PM은 접수자다), 구두 요청을 접수하지 않는 것, 모든 변경이 "긴급"인 것. 예상 소요 6분.}
\end{frame}

\begin{frame}{변경요청서의 구조 — 세 문서가 순서대로 만들어진다}
\begin{itemize}
\item \textbf{요청서}(항목 1\til{}3) — 식별 · 내용과 사유 · 분류와 긴급도
\item \textbf{영향 분석}(항목 4\til{}6) — 6축: 범위·일정·원가·품질·리스크·계약/편익
\item \textbf{의결서}(항목 7\til{}9) — 판정과 경로 · CCB 의결·조건·반대의견 · BL-0n
\item 결과는 \textbf{변경 로그} 한 줄(항목 10) — CR-01\til{}11 누적
\item 강조 칸 둘 — "현재 기준선 문장" 인용 · 대안 3개(요청대로/축소/이연·파라미터화/기각)
\end{itemize}
\keymsg{인용이 없는 변경은 무엇에서 무엇으로 바뀌는지 모르는 변경이다 — "원래 그랬다"와 "바뀐 것이다"가 충돌한다}
\note{교재 3.1, 워크북 §1.3 항목 1\til{}10. 요청서·영향 분석·의결서는 작성 시점도 작성자도 다르다 — 요청서는 접수 시, 영향 분석은 PMO·수행사 FP 산정 후, 의결서는 CCB 후. 그래서 한 양식이 아니라 세 문서가 이어 붙은 구조로 가르친다. 워크북 예시 CR-08의 항목 2를 보이며 "현재: 범위 기술서 v1.0 §2 ① … 헌장 §5 포함 범위에 임시 연동 없음 \ra{} 요청: … FP +90"의 인용 습관과, 대안 (a)\til{}(d) 넷 중 (b) A12 편입이 채택된 이유(코드 프리즈 위반 회피)를 읽는다. 긴급도의 근거는 날짜다 — CR-08 A12 착수 12-07, CR-09 FZ-10 10-28, CR-10 10-05 성과보고, CR-11 매핑 종료 10-23. 예상 소요 6분.}
\end{frame}

\begin{frame}{6축 영향 분석 — FP 하나로 변경의 크기를 잴 수 없다}
\begin{tbl}[\scriptsize]{L{1.7cm}L{4.6cm}L{4.6cm}}
축 & 무엇을 적는가 & 온담 CR-08 (+90 FP) \\
\midrule
범위 & WBS·RTM 델타, FP, \textbf{누적} FP & 1.1.7 산출물 +3, StRS-095, 누적 +210 \\
일정 & 활동·기간·\textbf{TF}, 임계경로 여부 & A12 20\ra{}25일, TF 10\ra{}5, 임계 불변 \\
원가 & 금액 + \textbf{자금원} & +0.495억, 컨틴전시 미배정 풀 \\
품질 & TC·커버리지·밀도 분모 & TC +12, 분모 12,400\ra{}12,490 \\
리스크 & 신규·변동 EMV, 허용범위 여유 & R-19 신규 0.60, 잔여 7.26(여유 0.74) \\
계약/편익 & 변경 대가·서면·B-ID·타 프로젝트 & 증액 0.495, 합의서 6호, P3 협의 09-22 \\
\end{tbl}
\keymsg{일정 축은 TF로, 원가 축은 자금원과 함께, 계약/편익 축은 타 프로젝트·운영 영향 칸과 함께 적는다}
\note{교재 3.2, 워크북 §1.3 항목 4\til{}6. 6축은 PRINCE2 Issues practice의 6개 성과 목표(범위·일정·원가·품질·리스크·편익)에서 왔고 온담은 편익 축에 계약을 붙였다 — 고정가 계약에서 변경은 곧 대가이기 때문이다. FP 하나로 재면 CR-10은 0이고 CR-11은 마이너스라 "영향 없음"으로 읽히지만, CR-10은 예외 보고 트리거를 바꾸고 CR-11은 계약 2건과 보존 요건을 건드린다. 각 축의 값은 Day2 기준선에서 인용한다 — A12 TF 10은 §3.5 CPM 표, 미배정 풀 5.60은 §4.5 항목 6, 잔여 EMV는 §5.5. 리스크 축의 "여유 0.74 — 경고"는 허용범위 8억에 7.26이 닿아가고 있다는 5교시의 복선이다. 예상 소요 8분.}
\end{frame}

\begin{frame}{변경 4건의 6축 히트맵 — 같은 자로 다른 변경을 잰다}
\fig[0.6]{../그림/PM_Day3/fig_3_2_six_axis_radar.pdf}
\figcap{그림 3-2. CR-08\til{}11의 6축 영향 히트맵과 처리 경로. CR-10은 FP 0인데 거버넌스 축이 진하고, CR-11은 FP $-$60인데 계약 축이 진하다}
\keymsg{FP 열만 보면 CR-10은 "영향 없음"이고 CR-11은 "좋은 소식"이다 — 6축이 그 착시를 걷어낸다}
\note{교재 3.2 그림 3-2. 네 건을 한 화면에 놓고 축별로 읽는다 — 범위 축은 CR-08이 가장 크고(+90), 일정 축은 CR-09가 임계경로 A10(이미 $-$10)에 +1일을 얹어 단독으로는 허용범위 내지만 ⑫ 회복 계획 전제이며, 계약/편익 축은 CR-09가 물류본부·3PL을 건드려 운영위 보고를 부른다. 히트맵 오른쪽 처리 경로 열이 다음 프레임의 규칙 6개로 이어진다. 수강생에게 "여러분 회사 변경 양식에 이 여섯 열 중 몇 개가 있는가"를 묻는다 — 대개 범위·원가 둘이다. 예상 소요 5분.}
\end{frame}

\begin{frame}{처리 경로 6규칙 — 허용범위는 크기의 기준, 헌장은 권한의 기준}
\begin{tbl}[\scriptsize]{L{0.5cm}L{5.6cm}L{4.8cm}}
\# & 조건 & 경로 \\
\midrule
① & 6축 모두 허용범위 내 · 헌장 §5·§12·§13 문언 불변 · 타 영향 없음 & \textbf{P1 PM 전결}(RACI 16행), 다음 CCB 사후 보고 \\
② & \textbf{헌장 문언 변경} 또는 어느 축 초과 & \textbf{CCB} 승인(17행); 수행사 FP 영향 시 \textbf{변경협의체} 합의 선행(18행) \\
③ & 관리예비비 필요 & \textbf{스폰서} 결재(20행) \\
④ & 타 프로젝트·운영 조직 영향 & 프로그램 \textbf{운영위} 보고 병행 \\
⑤ & 릴리스 이연(R5) & 헌장 §13 인용, CCB \\
⑥ & 승인된 변경 & CCB 후 5영업일 내 \textbf{BL 번호} 일괄 갱신 \\
\end{tbl}
\keymsg{CR-08(+90)·CR-11($-$60)은 허용범위 안이다 — 그러나 헌장 §5 문장을 바꾸므로 CCB다(규칙 ②)}
\note{교재 3.3, 워크북 §1.5 첫머리 "처리 경로 규칙(변경 관리 계획 v1.0 §3, G1 확정 [추가 설정])". 두 문이 다르다는 것이 이 교시의 핵심 문장이다 — 허용범위는 "얼마나 큰가"를 재고, 헌장은 "누가 서명했는가"를 잰다. 스폰서와 이사회가 서명한 문장을 PM이 전결로 바꾸면 다음 게이트에서 헌장과 기준선이 어긋난다. 순서도 중요하다 — 변경협의체(계약, 수행사와 FP 합의) \ra{} CCB(사내 의결) \ra{} 계약 서면. 변경협의체 없이 CCB가 FP를 정하면 수행사는 G4에서 다시 협상한다. 민간에는 공공의 과업심의위원회가 없으므로 이 협의체가 그 자리를 대신한다는 점, CR-07(+40, 헌장 문언 불변)은 규칙 ①로 전결·사후 보고된 대조 사례라는 점을 짚는다. 예상 소요 8분.}
\end{frame}

\begin{frame}{네 건이 같은 문을 통과한다 — 접수에서 BL-02까지}
\fig[0.6]{../그림/PM_Day3/fig_3_6_change_flow.pdf}
\figcap{그림 3-6. 접수 \ra{} 6축 \ra{} 경로 판정 \ra{} 변경협의체(09-30) \ra{} 운영위(10-08) \ra{} CCB(10-15) \ra{} 계약 서면(10-23) \ra{} BL-02(10-23). CR-07 전결은 대조}
\keymsg{변경협의체 \ra{} CCB \ra{} 계약 서면 \ra{} BL — 순서가 바뀌면 G4에서 대가를 다시 협상한다}
\note{교재 3.6 그림 3-6. 네 건의 날짜를 따라 읽는다 — 접수 09-16\til{}22(모두 FZ-10 접수 마감 10-14 전), 변경협의체 09-30 합의(CR-08·09·11 FP 확인, 이재현), 운영위 10-08(CR-09 물류·3PL 영향), CCB 제10차 10-15, 계약 서면 10-23(SI 변경 합의서 6호에 세 건 통합 +65 FP·+0.3575억, 용역 변경 계약 1호 $-$2.0), BL-02 발행 10-23. CR-07은 같은 그림의 옆길 — PM 전결 09-17, CCB 사후 보고. 규칙 ⑥ 5영업일이 10-15 \ra{} 10-23으로 지켜졌음을 확인시킨다. 예상 소요 5분.}
\end{frame}

\begin{frame}{CR-08 · CR-09 — 범위 추가와 요구 변경은 어떻게 다르게 읽히는가}
\begin{columns}[T]
\begin{column}{0.49\textwidth}
\subhead{CR-08 파일럿 임시 재고 동기화 (+90 FP)}
\begin{itemize}
\item 헌장 v1.1 수정 요청 1번 · 정미란 · 09-16
\item 대안 (b) A12 편입 — 코드 프리즈 회피
\item +0.495억 미배정 풀 · TF 10\ra{}5 · R-19 신규
\item 헌장 §5 문언 변경 \ra{} 변경협의체 \ra{} \textbf{CCB}
\end{itemize}
\end{column}
\begin{column}{0.49\textwidth}
\subhead{CR-09 가맹 발주 마감 20:00\ra{}22:00 (+35 FP)}
\begin{itemize}
\item 협의회 공문 · 서정필(S-10) · 09-17
\item 대안 (c) \textbf{파라미터화}, 초기 운용 20:00
\item +0.19억 · 임계 A10 +1일 · R-16 P 30\ra{}45
\item 물류·3PL 영향 \ra{} 규칙 ④ \textbf{운영위} + CCB
\end{itemize}
\end{column}
\end{columns}
\keymsg{CR-09가 CCB인 이유는 크기가 아니다 — P1 CCB는 이천센터 야간조 편성을 결정할 수 없다}
\note{워크북 §1.5 CR-08·CR-09 전문, 교재 3.6·해설 "왜 CR-09를 대안 (c)로". CR-08은 오류 정정에 가까운 범위 추가다 — WBS 1.1.7 카드가 "FP 0 — 배포 기록만"이었고, 2019년 스마트오더가 매장 운영 부하로 중단된 경로를 12개점에서 재현하지 않으려면 연동 2본(68 FP)과 배포·회수 도구(22 FP)가 필요하다. CR-09는 요구 자체는 타당하나(현행 웹 23:00보다 후퇴), 22:00 마감은 P1 화면이 아니라 21:00 피킹·폴백 야간 배치 21:00·03:00·대성 배송 출발 시각을 건드린다. 그래서 파라미터로 구현해 두고 운용 전환은 야간조 편성안 첨부 후 2027-01 CCB 재상정으로 조건을 달았다. 한동석(물류본부)은 CCB 위원이 아니고 파라미터 구현에 이의도 없었지만 서면 의견을 "반대의견" 칸에 실명 기재했다 — Day1 거버넌스의 CCB 규칙이자, 재상정 때 그의 의견이 출발점이 되게 하려는 장치다. 예상 소요 8분.}
\end{frame}

\begin{frame}{CR-10 · CR-11 — FP 0과 FP $-$60에도 6축을 다 돌린다}
\begin{columns}[T]
\begin{column}{0.49\textwidth}
\subhead{CR-10 원가 허용범위 정의 확정 (FP 0)}
\begin{itemize}
\item 헌장 §12 "+6\% (10.08)" \ra{} "컨틴전시 전액 10.2 (+7.0\%)"
\item 세 문서가 두 값 — 155.88 vs 156.0
\item 예외 보고 트리거 "EAC $>$ 156.0" 단일화
\item \textbf{CCB} + 스폰서 확인 + 운영위 보고
\end{itemize}
\end{column}
\begin{column}{0.49\textwidth}
\subhead{CR-11 전환 대상 5년 한정 ($-$60 FP)}
\begin{itemize}
\item 헌장 §5 ⑤ 7.4TB \ra{} 5년 이내 약 4.8TB
\item 재경팀장 보존 요건 확인서 09-22
\item SI \textbf{감액 $-$0.33억} · 용역 계약 11.4\ra{}9.4
\item A09 TF 5\ra{}15 · 계약 서면 2건 · CCB
\end{itemize}
\end{column}
\end{columns}
\keymsg{CR-10은 돈이 한 푼도 움직이지 않지만 예외 보고의 자를 바꾼다 — 감액 변경도 변경이다}
\note{워크북 §1.5 CR-10·CR-11 전문, 교재 해설 "왜 CR-11에도 6축을 다 돌렸는가". CR-10은 Day2 원가 기준선 항목 9에서 발견된 불일치의 확정이다 — 프로그램 규정 별표의 +6\%는 승인 예산 168.0을 분모로 한 프로그램 계층 규정이고, P1의 PM 허용범위는 "스폰서 보고 없이 쓸 수 있는 돈" = 컨틴전시 10.2다. 10.08로 두면 계획 원가 + 허용범위 = 155.88 $\neq$ 집행 한도 156.0이라 "허용범위 내인데 집행 한도 초과"가 논리적으로 가능해진다. 10-05 성과보고의 첫 EAC 기반 예외 판정이 이 정의를 쓰므로 긴급도가 있다. CR-11의 핵심은 원가가 아니라 보존 요건이다 — 확인서 없이 축소하면 세무조사 때 "5년 전 거래를 왜 못 보느냐"가 P1 책임이 된다. 원가 기준선 BL-01은 이미 9.4로 기반영되어 있었고 용역 계약만 11.4였으므로 이 변경은 계약을 기준선에 정합시키는 것이다. 관리예비비 복귀 조건 1건이 해제된다. 예상 소요 7분.}
\end{frame}

\begin{frame}{변경 로그 CR-01\til{}CR-11 — 누적 +185FP는 허용범위 ±372의 49.7\%}
\begin{tbl}[\scriptsize]{L{1.2cm}L{4.2cm}rrL{2.6cm}}
CR & 내용 & FP & 억 & 경로 \\ \midrule
03 & 인터페이스 코드 체계 보정 & +22 & +0.12 & PM 전결 \\
06 & 가맹 포털 알림 채널 추가 & +58 & +0.32 & CCB \\
07 & 인터페이스 미등재 2본 실체화(IS-03) & +40 & +0.22 & CCB \\
08 & 파일럿 임시 연동(헌장 §5 문언) & +90 & +0.495 & CCB · 조건 4 \\
09 & 22:00 마감 파라미터(조건부) & +35 & +0.19 & 협의체 \ra{} CCB \\
11 & 데이터 전환 5년 한정(감액) & $-$60 & $-$0.33 & CCB \\
\end{tbl}
\keymsg{감액도 변경이다 — CR-11이 로그에 없으면 "범위가 늘기만 했다"는 착시가 남고 변경 대가 협상에서 60FP를 잃는다}
\note{교재 3.6 "변경 로그", 워크북 §1.5 변경 로그 11건과 §5.5 변경 대가 산정. CR-01·02·04·05·10은 FP 0(절차·문언·거버넌스 변경)이라 표에서 뺐다 — 그러나 로그에는 있다는 것이 요점(변경 = FP가 아니라 기준선을 건드리는 모든 것). 누적 +185FP = 1.5\%, 허용범위 ±3\% = ±372의 49.7\%. 금액 +1.0175억(SI, 55만 원/FP)이 BL-02에서 PMB로 편입되고 미배정 풀 2.47이 그만큼 줄었다(3교시 앞 프레임 "네 건이 같은 문을 통과한다"의 결과). 학생에게 묻기 — "허용범위의 절반을 9개월에 썼다. 남은 8개월은?" 답은 정답이 없고, R-10(정본 판정 후 변경 요청)의 EMV 0.76이 그 질문의 정량판이다. 예상 소요 4분.}
\end{frame}

\begin{frame}{"Day1 문서는 고쳐진다" — 헌장 v1.1과 변경의 흐름}
\fig[0.48]{../그림/PM_Day3/fig_3_6_change_flow.pdf}
\figcap{그림 3-6. 접수 \ra{} 분류 \ra{} 6축 영향 \ra{} 경로 판정 \ra{} 의결 \ra{} 기준선 갱신 BL-0n \ra{} 문서 역전파(헌장·RTM·WBS 사전·등록부·PV)}
\keymsg{헌장을 고치는 것은 실패가 아니다 — 고치지 않은 헌장이 실패다. 단, 고친 흔적(v1.1, CR 번호, 승인자)이 없으면 그것은 변경이 아니라 변조다}
\note{교재 3.7, 워크북 §1.6 해설과 §6. Day2 §6이 남긴 헌장 v1.1 수정 요청 10건 중 오늘 CCB에 오른 것이 CR-08·10·11(3건)이고, 나머지는 BL-02 발행 시 문서 갱신으로 흡수됐다. 역전파 순서 — 헌장 §5 문언(CR-08) \ra{} 범위 기술서·WBS v1.1 \ra{} RTM v1.2 \ra{} PV(2026-12 +0.33·2027-01 +0.69) \ra{} 등록부 v1.2. 이 순서가 거꾸로 되면(PV부터 고치고 헌장은 나중에) 감리가 "기준선 근거 없음"으로 잡는다 — 감리 1번 지적이 정확히 그것이었다. Day1 슬라이드 "헌장은 권한 문서다"를 다시 띄우고 "권한 문서를 누가 고칠 수 있는가"로 닫는다(스폰서 서명 = 정미란). 예상 소요 4분.}
\end{frame}

\begin{frame}{3교시 핵심 정리}
\begin{enumerate}
\item 변경은 사건이 아니라 절차 — 접수·6축·경로·의결·갱신의 다섯 단계 중 하나라도 건너뛰면 변조다
\item 6축 영향 분석은 FP 하나로 대체되지 않는다 — CR-10은 FP 0인데 원가 축 정의를 바꿨다
\item 처리 경로 6규칙 — 허용범위는 크기의 기준, 헌장 문언은 권한의 기준, 운영 영향은 소관의 기준
\item BL-02: FP 12,585 · PMB 146.82 · 집행 한도 156.00 불변 — 기준선 갱신은 승인된 변경으로만
\item 로그 누적 +185FP(49.7\%)와 감액 $-$60 — 변경 로그는 협상의 근거이자 허용범위의 잔고
\end{enumerate}
\keymsg{4교시 실습 — CR-09를 직접 6축으로 재고 의결서를 쓴 뒤, 9/30 세 숫자에서 열두 값을 손으로 만든다}
\note{교재 제3장 핵심 정리 1\til{}5. 생각해 볼 질문(제3장 3 — "여러분 회사에서 마지막으로 헌장이 개정된 것은 언제인가, 개정된 적이 없다면 그것은 무엇을 뜻하는가"). 점심 전 마지막 이론 교시이므로 4교시 준비물(워크북 §1.4·§2.4 빈 템플릿, 계산기)을 안내한다. 예상 소요 3분.}
\end{frame}

% ===== 4교시 (90분) — 실습 ① 변경요청서·영향 분석·CCB 의결서 + EVM·ES 손계산 =====
\section{4교시 (90분) — 실습 ① 변경요청서·영향 분석·CCB 의결서 + EVM·ES 손계산}

% =====================================================================
% 4교시 — 실습 ①
% =====================================================================
\begin{frame}{이 교시에서 하는 것 — 오전 이론 둘을 손으로 쓴다}
\begin{itemize}
\item \textbf{⑪ 변경요청서} — 변경 4건 중 \textbf{CR-09(가맹 발주 마감 20:00 \ra{} 22:00)}를 6축으로 재고 경로를 판정한다
\item \textbf{⑫ 성과보고} — 9/30 누적 PV 77.83 · EV 73.43 · AC 71.13에서 지수·EAC·ES를 \textbf{손으로} 계산한다
\item 두 문서는 서로의 입력 — CR-09의 FP +35가 EAC4의 ETC 행에, A10 $-$10이 CR-09 일정 축에
\item 도구는 종이 캔버스 · 계산기 · 워크북 §1.4·§2.4 빈 템플릿 — 노트북은 검산용
\item 시간 — 안내 10분 / ⑪ 35분 / ⑫ 35분 / 짝 교환 검산 10분
\end{itemize}
\keymsg{정답은 워크북 §1.5·§2.5에 있다 — 보지 말고 쓰고, 다 쓴 뒤에 대조하라}
\note{워크북 §1·§2 실습. Day2 4교시와 같은 진행 — 과제 지시 두 프레임, 작성 요령 네 프레임, 흔한 오류 한 프레임, 예시본 대조, 정리. 오전에 배운 3교시(6축·처리 경로 규칙 ①\til{}⑥)와 1·2교시(편차·지수·EAC 4공식·Earned Schedule)가 이 90분의 재료다. 두 과제를 하나의 교시에 묶은 이유를 먼저 말한다 — 변경요청서의 원가 축(+0.19억)은 성과보고의 상향식 ETC "SI 승인·예정 변경 대가 1.02" 행의 일부이고, 성과보고의 A10 TF 0·$-$10은 변경요청서 일정 축의 전제다. 한쪽만 쓰면 다른 쪽이 틀린다. 시간 배분을 화면에 두고 35분 단위로 종을 울린다. 예상 소요 3분.}
\end{frame}

\begin{frame}{과제 ⑪ — CR-09 가맹 발주 마감 20:00 \ra{} 22:00, 항목 1\til{}8을 직접 쓴다}
\begin{itemize}
\item 접수 — 2026-09-17 협의회 공문 제2026-14호, 요청자 S-10 서정필 협의회장, FZ-10(10-28) 전
\item 기준선 문장 — RTM StRS-073 "접수\til{}출고 확정 12h" \ra{} FRN-FLOW-01 §2 "마감 20:00, 피킹 21:00, 확정 익일 08:00"
\item 대안 넷 — (a) 22:00·야간조 연장 (b) 22:00·14h (c) \textbf{파라미터화, 초기 20:00} (d) 기각
\item 쓸 것 — 6축 영향(항목 4\til{}6) · 허용범위 판정과 처리 경로(항목 7) · CCB 의결서(항목 8)
\item 주어진 것 — 수행사 FP 산정 +35(파라미터 12 + 배치 15 + SAP 시각 8), A10 TF 0 · 현재 $-$10
\end{itemize}
\keymsg{묻는 것은 하나 — 이 변경을 P1 CCB가 결정할 수 있는가, 결정할 수 없는 부분은 어디로 보내는가}
\note{워크북 §1.4 빈 템플릿에 쓴다. CR-08(스폰서 범위 추가)·CR-10(거버넌스)·CR-11(범위 축소)은 예시본으로 읽기만 하고, CR-09만 학생이 쓴다 — 넷 중 유일하게 "P1 권한 밖의 것"이 섞여 있어 처리 경로 판정이 연습이 되기 때문이다. 항목 1\til{}3은 위 정보를 옮겨 적으면 되므로 5분 안에 끝내게 하고, 35분의 대부분을 항목 4\til{}8에 쓰게 한다. 학생이 자주 묻는 것 — "대안을 하나 골라야 하나?" 그렇다. CCB 의결서는 대안 하나에 조건을 붙여 의결하는 문서이고, 대안 넷을 나열만 한 의결서는 의결이 아니다. 추정치가 필요한 곳(R-16 확률 변화, TC 건수)은 "[추정]"을 붙여 적게 한다 — 숫자를 비워 두는 것보다 근거 있는 추정이 낫다. 예상 소요 4분.}
\end{frame}

\begin{frame}{과제 ⑫ — 9/30 누적 세 숫자에서 열두 값을 손으로}
\begin{tbl}[\scriptsize]{L{2.6cm}L{5.4cm}L{3.2cm}}
단계 & 계산할 것 & 입력 \\ \midrule
편차·지수 & SV · CV · SPI · CPI & PV 77.83 · EV 73.43 · AC 71.13 \\
예측 & EAC1 · EAC2 · EAC3 · EAC4 & BAC 145.80 · 상향식 ETC 79.72 \\
목표 효율 & TCPI(BAC) · TCPI(집행 한도 156.00) & BAC · 156.00 \\
Earned Schedule & ES · SV(t) · SPI(t) · IEAC(t) & PV(8월) 71.74 · AT 9 · PD 17 \\
\end{tbl}
\begin{itemize}
\item 소수 셋째 자리까지, 산식을 값 옆에 반드시 쓴다 — 산식 없는 값은 검산 불가
\item 본값을 먼저 끝내고, 시간이 남으면 조정(시차 1.96 · 라이선스 4.70)을 한 줄 더
\end{itemize}
\keymsg{계산기는 써도 된다 — 공식을 외우지 말고, 각 공식의 전제를 값 옆에 한 단어로 적어라}
\note{워크북 §2.4 빈 템플릿 항목 4\til{}6. 월별 PV·EV·AC 표(항목 3)는 이미 채워진 것을 나눠 주고 누적 행만 쓰게 한다 — 이 교시의 목적은 표 작성이 아니라 지수와 예측의 손계산이다. 상향식 ETC 79.72는 항목별 근거표(SI 잔여 38.51, 변경 대가 1.02, 상용SW 5.90, 인프라 7.92, 단말 8.90, 전환 5.10, 시험·감리 7.40, PMO 4.80, 대기 0.18)를 주고 합만 확인하게 한다. 학생이 EAC1이 BAC보다 작게 나오면 당황하는데, 그것이 정답이다 — 141.23이 나오는 이유를 "고정가에서 PMB 아래 EAC는 불가능"이라고 판단하는 것이 2교시의 착시 논의를 손으로 확인하는 지점이다. ES는 PV(8월) 71.74와 PV(9월) 77.83 사이에서 EV 73.43이 어디에 있는지 보간하는 것 — 그림 2-4를 벽에 띄워 둔다. 예상 소요 4분.}
\end{frame}

\begin{frame}{작성 요령 ① — 6축의 각 축에 무엇을 적는가}
\begin{tbl}[\scriptsize]{L{1.4cm}L{5.6cm}L{4.0cm}}
축 & 적는 것 & CR-09에서 \\ \midrule
범위 & WBS 번호 · RTM 변경 이력 · SRS 건수 · FP 델타와 \textbf{누적} & 1.4.2 · 1.5.4 · SRS 4건 · +35 (누적 +245) \\
일정 & 어느 활동의 \textbf{TF}가 얼마 줄었나 — 일수가 아니라 여유 & A10 TF 0 · $-$10에 +1영업일 \\
원가 & FP $\times$ 단가 · 자금원(컨틴전시 풀 / 관리예비비) & 35 $\times$ 0.0055 = +0.19 · 미배정 풀 \\
품질 & TC 재작성·추가 건수 · UAT 수정 · 커버리지 & TC 3건 재작성 · +8 · UAT-FRN-039 \\
리스크 & 어느 R-의 P·EMV가 움직이나 · 잔여 EMV 합 순증 & R-16 0.24 \ra{} 0.36 · R-02 하락 · 순증 +0.05 \\
계약/편익 & 변경 대가 · 편익 ID 유지 여부 · \textbf{타 프로젝트·운영 영향} & +0.19 · B-08 12h 유지 · 물류본부·3PL \\
\end{tbl}
\keymsg{여섯 칸은 모두 "무엇에서 무엇으로"다 — 기준선 문장을 인용하지 않은 축은 빈 칸과 같다}
\note{워크북 §1.3 항목 4\til{}6 작성 가이드. 학생이 가장 잘못 쓰는 축은 일정과 계약/편익이다. 일정 축에 "+1일"이라고 쓰면 아무 뜻이 없다 — 어느 활동인지, 그 활동의 TF가 얼마인지, 임계경로인지가 있어야 허용범위 판정(컷오버 창 0)이 된다. CR-09는 A10이 이미 $-$10이므로 "+1영업일이 허용범위 내"라는 판정은 성과보고의 회복 계획을 전제로 해서만 성립한다 — 그 전제를 판정 칸에 적는 것이 정답이다. 계약/편익 축의 "타 프로젝트·운영 영향" 칸은 CR-09의 핵심이다. 22:00 마감은 P1 화면의 변경이 아니라 이천센터 야간 작업과 대성로지스 배송 출발 시각의 변경이며, 그 칸이 있어야 처리 경로 규칙 ④(운영위 보고 병행)가 발동한다. 범위 축의 누적 +245는 CR-08까지의 +210에 +35를 더한 것 — 허용범위 ±372의 66\%다. 예상 소요 5분.}
\end{frame}

\begin{frame}{작성 요령 ② — 처리 경로 판정과 CCB 의결서의 네 칸}
\begin{itemize}
\item 판정 순서 — 6축 허용범위 \ra{} 헌장 문언 변경 여부 \ra{} 타 프로젝트·운영 영향 \ra{} 관리예비비 필요
\item CR-09 — 6축 모두 내 · 헌장 문언 불변 · \textbf{운영 영향 예} \ra{} 규칙 ④ · PM 전결 불가
\item 경로 — 변경협의체(09-30 합의) \ra{} CCB(10-15) + 프로그램 운영위 보고(10-08)
\item 의결서 네 칸 — \textbf{의결}(대안 하나) · \textbf{조건}(누가 언제까지) · \textbf{반대의견}(서면 그대로) · \textbf{미이행 시}
\item 설계 변경과 운영 변경을 \textbf{분리}한다 — 파라미터(P1 권한)는 승인, 22:00 운용(물류 권한)은 재상정
\end{itemize}
\keymsg{CCB가 자기 권한 밖의 것을 결정하면 집행되지 않고, 미루면 요구가 방치된다 — 분리해서 의결하라}
\note{워크북 §1.3 항목 7·8과 §1.5 처리 경로 규칙 ①\til{}⑥. 판정 순서를 화면 순서대로 고정시킨다 — 학생들은 흔히 "FP +35는 작으니 PM 전결"로 끝내는데, 허용범위 내라는 것은 규칙 ①의 첫 조건일 뿐이고 규칙 ④(운영 영향)가 걸리면 CCB다. 의결서의 정답 골격 — 조건부 승인 · 대안 (c) · FP +35 · 0.19억 · 초기 운용 20:00. 조건 ① 22:00 전환은 물류본부 야간조 편성안과 3PL 배송 출발 시각 합의를 첨부해 2027-01 CCB에 재상정, ② 협의회 회신 10-23, ③ B-08 12h 유지·14h 불채택, ④ R-16 잔여 EMV 갱신. 미이행 시 — 2027-01 CCB 미상정이면 20:00 운용 확정. 반대의견 칸에 한동석의 서면("야간조 편성은 P2 안정화 이후, 파라미터 구현에는 이의 없음")을 그대로 적는 것이 의결서를 나중에 읽는 사람을 위한 최소한이다. 이해관계자 퇴장(이재현, 변경 대가 확인 후 의결 시)도 기록 항목이다. 예상 소요 5분.}
\end{frame}

\begin{frame}{작성 요령 ③ — EVM 손계산의 순서: 세 숫자에서 열두 값으로}
\begin{itemize}
\item 입력 셋 — PV 77.83 · EV 73.43 · AC 71.13 (누적, 2026-09-30) · BAC 145.80
\item 편차 둘 — SV = EV $-$ PV = $-$4.40 · CV = EV $-$ AC = +2.30 (부호를 먼저 적는다)
\item 지수 둘 — SPI = EV/PV = 0.943 · CPI = EV/AC = 1.032 (소수 셋째 자리)
\item EAC 넷 — BAC/CPI = 141.23 · AC+(BAC$-$EV) = 143.50 · AC+(BAC$-$EV)/(CPI×SPI) = 145.43 · AC+상향식 ETC 79.72 = 150.86
\item TCPI 둘 + VAC — TCPI(BAC) = (BAC$-$EV)/(BAC$-$AC) = 0.969 · TCPI(156.0) = 0.853 · VAC = BAC $-$ EAC (넷 다)
\end{itemize}
\keymsg{순서를 바꾸지 마라 — EAC를 먼저 쓰고 CPI를 맞추는 학생이 매 기수 있다. 공식은 전제의 이름이지 결과의 이름이 아니다}
\note{워크북 §2.3 항목 3\til{}5, §2.6 검산식(src/day3\_evm.py). 판서 순서를 이 다섯 줄 그대로 고정한다. 학생이 가장 자주 틀리는 자리 셋 — (1) SV 부호(PV$-$EV로 계산해 +4.40), (2) EAC3의 분모(CPI×SPI를 CPI+SPI로), (3) TCPI 분모에 BAC 대신 EAC를 쓰고도 "TCPI(BAC)"라고 적는 것. 상향식 ETC 79.72의 내역(SI 잔여 38.51 + 변경 1.02 + 대기 0.18 + 상용SW 5.90 + 인프라 7.92 + 단말 8.90 + 전환 5.10 + 시험·감리 7.40 + PMO 4.80)은 워크북 표를 옆에 두고 옮겨 적게 한다 — 이것이 "공식이 아니라 근거"라는 2교시 명제의 실습판. 착시 조정(AC 75.83 \ra{} CPI 0.968)은 심화 과제. 예상 소요 15분(학생 작업 포함).}
\end{frame}

\begin{frame}{작성 요령 ④ — Earned Schedule 읽기: EV가 PV 곡선에서 "언제"인가}
\begin{itemize}
\item ES = 8 + (EV $-$ PV$_8$) / (PV$_9$ $-$ PV$_8$) = 8 + (73.43 $-$ 71.74) / (77.83 $-$ 71.74) = \textbf{8.28개월}
\item AT = 9 (데이터 일자 2026-09-30 = 착수 후 9개월)
\item SV(t) = ES $-$ AT = $-$0.72개월 (≈ $-$15영업일) · SPI(t) = ES/AT = \textbf{0.920}
\item IEAC(t) = PD / SPI(t) = 17 / 0.920 = \textbf{18.48개월} \ra{} 추세 종료 2027-07-12 (허용 G4+15 = 06-18)
\item 대조 — SPI 0.943(돈) vs SPI(t) 0.920(시간): 7월부터 벌어지고 9월 차이 $-$0.023
\end{itemize}
\keymsg{ES는 "EV 73.43을 PV 곡선이 언제 지났는가"를 자로 재는 것 — 그림 2-4의 빨간 점을 손으로 찾으면 끝난다}
\note{워크북 §2.3 항목 6, 그림 2-4. 계산보다 읽기가 먼저 — 그림 2-4를 배포하고 EV 73.43을 PV 곡선에서 찾아 x축을 읽게 한다(8월과 9월 사이, 8월에 더 가까움). 그다음 보간식으로 8.28을 확인한다. PV$_8$ = 71.74·PV$_9$ = 77.83은 Day2 §4 월별 PV 표의 정본. IEAC(t)를 날짜로 옮길 때 17개월 PD의 기준일(2026-01-05)에서 18.48개월을 더하면 2027-07-12 — 허용범위 종료 06-18보다 늦으므로 일정 축 예외(EX-01)의 정량 근거가 된다. SPI와 SPI(t)의 차이가 왜 7월부터 벌어지는지(PV 곡선이 5\til{}6월 조달 봉우리 뒤로 완만해져 같은 EV 부족이 더 긴 시간으로 환산됨)는 심화 질문. 예상 소요 12분(학생 작업 포함).}
\end{frame}

\begin{frame}{흔한 오류 — 이 실습에서 매 기수 나오는 다섯}
\begin{itemize}
\item 구두 합의를 변경으로 쓰지 않는다 — "서정필 회장과 이야기됐다"는 CR-09의 접수 근거가 아니라 접수 사유다
\item SPI를 일수로 읽는다 — 0.943 × 9개월 ≠ 지연 일수. PV 곡선에 봉우리가 있으면 틀린다(ES를 쓰라)
\item 월별 지수와 누적 지수를 섞는다 — 4월 누적 SPI 0.784와 4월 월별 EV/PV는 다른 숫자
\item 라이선스 이연을 원가 절감으로 읽는다 — CV +2.30의 +4.70은 아직 안 낸 돈이다
\item 6축 중 "편익" 축을 비운다 — CR-09는 편익 B-08(리드타임 12h)에 직접 걸리는 변경이다
\end{itemize}
\keymsg{다섯 오류의 공통점 — 숫자를 계산은 했는데 그 숫자가 무엇의 이름인지 적지 않았다}
\note{워크북 §1.7·§2.7 점검 질문과 Day1\til{}Day2 실습에서 반복 관찰된 오류를 모았다. 각 조 산출물을 걷기 전에 이 다섯을 화면에 띄우고 자기 것을 스스로 고치게 한다(10분). 특히 4번 — 5월 CPI 1.073 점프를 "원가 효율 개선"으로 적은 조가 있으면 그 조의 성과보고서를 예시로(익명) 띄워 착시 1을 다시 설명한다. 예상 소요 8분.}
\end{frame}

\begin{frame}{예시본 참조 · 시간 배분 · 심화 과제}
\begin{itemize}
\item 예시본 — 워크북 §1.5(CR-09 완성본 + CCB 10-15 의결서) · §2.5(성과보고 제9호, Highlight Report 1장)
\item 인쇄용 — \textsf{../템플릿/PM\_Day3/11\_변경요청서\_완성예시.pdf} · \textsf{12\_성과보고\_완성예시.pdf} (빈 양식은 \textsf{\_빈템플릿.pdf})
\item 시간 — ⑪ 변경요청서·6축·의결서 40분 / ⑫ EVM·ES 손계산 35분 / 자기 점검·대조 15분
\item 자기 점검 — 워크북 §1.7·§2.7 점검 질문 각 3개. 예시본과 값이 다르면 \textbf{어느 전제가 달랐는지}를 적는다
\item 심화(3\til{}7년차) — 자기 프로젝트의 최근 변경 1건을 6축으로 다시 재고, 마지막 성과보고에 SPI(t)를 추가해 본다
\end{itemize}
\keymsg{예시본은 정답이 아니라 전제가 적힌 답이다 — 다른 값이 나왔으면 전제를 비교하라}
\note{Day2 4교시와 같은 운영. 예시본 PDF는 실습 시작 시 배포하지 않고 35분 뒤(⑫ 착수 시점)에 ⑪ 예시본만 먼저 배포, ⑫ 예시본은 손계산이 끝난 뒤. 값이 다른 조가 나오면 "전제"를 찾게 하는 것이 이 실습의 목적이다 — EAC3에서 CPI×SPI에 조정 CPI 0.968을 쓴 조는 155.04(경고 대역)가 나오고, 그것은 틀린 것이 아니라 다른 전제다. 예상 소요 90분 전체.}
\end{frame}

\begin{frame}{4교시 핵심 정리}
\begin{enumerate}
\item CR-09는 FP +35라 작아 보이지만 운영 영향(규칙 ④)으로 CCB — 크기와 권한은 다른 축이다
\item 의결서 네 칸(의결·조건·반대의견·미이행 시) 중 반대의견 칸이 나중에 읽는 사람을 위한 최소한
\item 세 숫자(PV·EV·AC)에서 열두 값 — 순서가 곧 전제. EAC 넷은 네 개의 미래다
\item ES 8.28은 계산이 아니라 읽기 — 그림에서 찾고 식으로 확인한다
\item 다섯 오류의 공통점 — 계산은 했는데 이름을 안 적었다
\end{enumerate}
\keymsg{5교시는 나머지 셋(⑬ 이슈·리스크 / ⑭ 품질·감리 / ⑮ 조달·계약) — 숫자가 아니라 기록이 무기가 되는 영역}
\note{조별 산출물 회수. 성과보고서 1장(Highlight Report)은 6교시 시작 전까지 벽에 붙여 두고 갤러리워크로 서로 대조하게 한다 — RAG 색이 다른 조가 있으면 그것이 5교시 "RAG의 정치" 도입 재료다. 예상 소요 3분.}
\end{frame}

% ===== 5교시 (75분) — 품질 · 이슈 · 리스크 · 조달 통제 =====
\section{5교시 (75분) — 품질 · 이슈 · 리스크 · 조달 통제}

% =====================================================================
% 5교시 — 품질 · 이슈 · 리스크 · 조달 통제 (교재 제4장·제5장, 워크북 §3~§5)
% =====================================================================
\begin{frame}{5교시에서 배우는 것 — 나머지 세 산출물(⑬·⑭·⑮)}
\begin{itemize}
\item 품질 — 결함밀도의 분모·분자를 누가 정했는가, 곡선의 평탄은 수렴인가 정지인가
\item 감리 — 등급보다 \textbf{원인 귀속} 칸이 나중의 책임 귀속을 정한다
\item 이슈·리스크 — 이슈에는 기한, 리스크에는 트리거. 잔여 EMV 7.26 vs 허용 8.0
\item 보고 — RAG의 정치와 mum effect. 나쁜 소식은 왜 올라오지 않는가
\item 계약 — 동시기록(검수 보류 서면 회신)·동시 지연 분해·수행사 성과 기록
\end{itemize}
\keymsg{네 번째 산출물부터는 숫자보다 "누가 원인이고 누가 결정하는가"를 적는 칸이 무겁다}
\note{교재 제4장(품질·이슈·리스크)과 제5장(사람·커뮤니케이션·조달)을 75분에 다룬다. 워크북 §3(⑬ 이슈 로그·리스크 등록부 갱신), §4(⑭ 품질 통제 기록), §5(⑮ 조달·계약 변경 기록)의 완성 예시본이 정본이다. 1\til{}4교시가 "편차를 재는 법"이었다면 5교시는 "편차의 원인을 누구에게 귀속시키고, 그 기록이 검수·변경협의체·지체상금에서 어떻게 쓰이는가"다. 세 산출물은 서로를 인용한다 — 감리 지적 7번(⑭)이 이슈 IS-07(⑬)이고 R-20(⑬)이며 검수 보류 회신(⑮)의 4항이다. 수강생에게 워크북 §4.5 항목 6 대응 관리대장과 §5.5 항목 4 서면 회신을 펴 두게 한다. 예상 소요 4분.}
\end{frame}

\begin{frame}{QA와 QC, 그리고 결함밀도 정의의 정치 — 분모와 분자}
\begin{itemize}
\item QA는 절차가 지켜졌는가(감리·검토회), QC는 산출물이 기준을 넘었는가(결함·밀도)
\item 분자 — 스프린트 시험 결함만(설계 검토 제외), 요구 변경·해석 차이는 결함이 아니라 CR
\item 분모 — \textbf{인수 FP 누적} 7,020. 계획 FP(7,440)로 나누면 0.125로 "이내"가 된다
\end{itemize}
\begin{tbl}[\scriptsize]{L{3.2cm}L{3.6cm}rL{2.3cm}L{2.6cm}}
계층 & 산식 & 값 & 기준 & 판정\\\midrule
작업패키지 & 930 $\div$ 7,020 FP & 0.132 & 0.13 (수행사 PM) & \textbf{초과 +1.5\%} \ra{} 1단계\\
설계 검토 포함(참고) & 1,140 $\div$ 7,020 & 0.162 & — & 분모 부정합\\
프로젝트(전망) & 0.132 $\times$ 유출 35\% + UAT 0.04 & 약 0.09 & 0.4 + 심1·2 미해결 0 & 이내 전망\\
\end{tbl}
\keymsg{0.132 > 0.13 — 1.5\% 초과라도 허용범위 밖이면 에스컬레이션한다. 반올림으로 "이내"를 만들지 않는다}
\note{교재 4.1, 워크북 §4.5 항목 1·4. 결함밀도는 정의가 곧 정치다 — 분자에 설계 검토 결함을 넣으면 0.162, 분모를 계획 FP로 바꾸면 0.125가 되고, 세 숫자 중 어느 것이 "기준 0.13"과 비교되는지는 Day2 격자에서 정했다(스프린트 시험 결함 ÷ 인수 FP). 그래서 문서 항목 1에 측정 정의를 먼저 적는다. 0.132는 0.13을 1.5\% 넘었을 뿐이지만 예외에 의한 관리에서 허용범위는 허용범위다 — 이재현이 인지 09-25 다음 영업일 09-28에 P1 PM에 보고했고, P1 PM은 3영업일 내(10-01)에 결함 해소 전담 2명 배치·backlog 10-31 60 이하·시험 노력 주간 보고 추가를 결정했다. 프로젝트 계층 0.4는 넉넉하고 실질 제약은 "심각도 1·2 미해결 0"임을 짚는다. 예상 소요 6분.}
\end{frame}

\begin{frame}{결함 누적곡선 — 평탄에는 두 가지 뜻이 있다}
\fig[0.58]{../그림/PM_Day3/fig_4_2_defect_curves.pdf}
\figcap{그림 4-2. 34주 누적 발견 1,140 · 해소 1,024 · 미해결 116. 3월 말 평탄(①)은 수렴, 8월 평탄(③)은 주당 TC 210 \ra{} 95로 시험이 멈춘 것 — 곡선만 보면 둘은 같아 보인다}
\keymsg{발견 곡선이 평평해지면 먼저 시험 노력을 본다. 노력 없이 줄어든 결함은 품질이 아니라 정지다}
\note{교재 4.2, 워크북 §4.5 항목 2·3. 네 구간을 읽는다 — ① 설계 검토 2\til{}3월 주 18\ra{}34\ra{}15, 검토 대상 47종이 S5까지 대부분 끝나 노력 감소 없이 발견이 줄었으므로 수렴형. ② S7\til{}S15 인수 FP에 비례해 상승, 6월 밀도 0.16 최고(R-05 지연 후 S11·S12 재작업). ③ 8월 12\ra{}9\ra{}11 — 시험 인력 6명 중 4명이 규제 검증(A07 142개 항목)으로 이동해 주당 TC가 210에서 95로 떨어졌다. 8월 밀도 0.129는 착시다. ④ 9월 재상승 27\ra{}52\ra{}28, 8월 미시험분이 S17에 몰렸다. 지금 봐야 할 것은 발견이 아니라 backlog다 — 6월부터 해소가 발견을 못 따라가 116건이 쌓였고 R3 개발자 시간을 잠식한다(A10 $-$10의 원인 중 하나). 10월부터 시험 노력(TC 실행)을 결함표 열로 추가하는 이유다. 예상 소요 6분.}
\end{frame}

\begin{frame}{심각도 2 미해결 7건 — 건수보다 나이를 본다}
\begin{tbl}[\scriptsize]{lL{3.6cm}rlL{3.4cm}}
ID & 모듈 (WBS) & 나이 & 기한 & 연결\\\midrule
D-0966 & API GW 인증 (1.5.2) & 16 & 10-16 & IS-08 보안 High, 감리 6번\\
D-0871 & 3PL 폴백 전환 (1.3.3) & 12 & 10-09 & R-16, TC-INV-037\\
D-0902 & POS 반품 적립 취소 (1.1.2) & 10 & 10-09 & 멤버십 StRS\\
D-0915 & 개인정보 파기 로그 (1.5.1) & 8 & 10-16 & 감리 6번, SRS-708\\
D-0958 & MDM 중복 검출 (1.5.3) & 4 & 10-23 & R-10, 정본 판정 전 필수\\
\end{tbl}
\begin{itemize}
\item 미해결 116 = 심1 \textbf{0} / 심2 \textbf{7} / 심3 58 / 심4 51 — 인수에는 지장 없으나 시간을 잠식
\item 나이(영업일) 중앙값 9, 심2 최장 16 — "심2 나이 10일 초과"를 주간 PM 회의 고정 안건으로
\item 심각도 기준서 v1.0(09-28)은 감리 5번 대응 — 116건 재분류 10-09
\end{itemize}
\keymsg{통합시험 착수(12-07) 조건은 심2 미해결 0·심3 30 이하 — G3의 "심1·2 = 0"은 지금부터 만든다}
\note{교재 4.2 "심각도와 나이", 워크북 §4.5 항목 5(7건 중 5건만 표시 — D-0931·D-0944는 워크북 참조). 결함 건수는 인수 FP에 비례해 늘기 마련이므로 절대 건수는 정보가 적다. 정보가 있는 것은 심각도별 미해결 수와 그 나이다 — D-0966은 보안 High(VA-1-01)와 같은 건이고 16영업일째 열려 있으며, D-0915는 감리 6번(파기 증적)과 원인이 같다. 결함 하나가 이슈(IS-08)·감리 지적·리스크(R-16·R-10)에 동시에 연결된다는 점을 보여 준다 — 그래서 연결 열이 있다. 우선순위(P1\til{}P3)는 심각도와 별도 열이라는 것도 짚는다. 예상 소요 5분.}
\end{frame}

\begin{frame}{감리 지적 8건 — 등급보다 원인 귀속 칸이 무겁다}
\fig[0.56]{../그림/PM_Day3/fig_4_3_audit_attribution.pdf}
\figcap{그림 4-3. 감리 기재 수행사 6·발주사 2 \ra{} 발주사 판정 수행사 3·발주사 4·공동 1. AF-1-03(P2 규격 확인 서명)·AF-1-04(가맹 SW 배포 범위)는 WBS 담당 열·RACI로 발주사 귀속을 입증해 09-29 서면 이의}
\keymsg{지적 사실은 인정하고 귀속만 다툰다 — 그 두 글자가 R2 검수·대기 비용·변경협의체의 근거가 된다}
\note{교재 4.3, 워크북 §4.5 항목 6과 이의 제기 서면 완성본. 감리 지적 대응 관리대장의 핵심 열은 등급(중대·보통·경미)이 아니라 "감리 기재 귀속"과 "발주사 판정 귀속"의 두 열이다. AF-1-03은 규격이 WBS 1.3.1 발주 담당 산출물이고 P2 전달은 RACI 9행 A = 프로그램 PMO장인데, 설비 벤더의 05-29 변경 요청에 프로그램 PMO가 09-30 현재 회신하지 않았다(회신 요청 06-05·07-10·08-21, 수행사 미결 통지 06-08·08-24 기록). AF-1-04는 1.1.6 대응안 결정이 오정근 A이며 05-08 "부속서 협상 후로 유보" 회신이 있다. 이의 서면의 구조 — ① 지적 사실 인정 ② 귀속 재분류 요청과 근거(WBS 담당 열·RACI·기록) ③ 발주사 조치 계획(10-16·11-27) ④ 요청(정본 정정, 10-08 회신, 향후 귀속 판정 기준 합의). 등급을 다투지 않는 이유 — 등급은 감리의 권한이고, 귀속은 사실의 문제다. 예상 소요 6분.}
\end{frame}

\begin{frame}{이슈 로그 11건 — 이슈에는 기한이, 리스크에는 트리거가 있다}
\begin{tbl}[\scriptsize]{lL{4.2cm}lL{2.2cm}l}
ID & 제목 & 심 & 에스컬 & 상태\\\midrule
IS-01 & R-04 실현 — 2026 배분 $-$8.0 & 1 & 2단계(스폰서) & 종결 03-13\\
IS-02 & R-05 실현 — POS 화면 승인 12영업일 & 1 & 병행(변경협의체) & 종결 07-15\\
IS-05 & R-03 트리거 — 3PL 폴백 확정 & 2 & 병행(임원\ra{}스폰서) & 종결 08-21\\
IS-07 & 감리 8건 — 3·4번 귀속 이의 & 1 & 2단계(회신 10-08) & 진행\\
IS-08 & 보안진단 High 2건 & 1 & 병행(최영주) & 진행\\
IS-11 & 수행사 인터페이스 PL 교체 & 2 & 1단계 & 종결 08-03\\
\end{tbl}
\begin{itemize}
\item 출처 — 리스크 전환 3 / 신규 발견 2 / 감리·진단 2 / 외부 3 / 수행사 1 (종결 5 · 진행 6)
\item 등록부에 없던 원인 5건 중 3건(IS-04·08·11)은 프리모템 원문에 있었다 \ra{} Day4 교훈 ⑲
\item 트리거 발동은 신호(기록·집행 준비), 실현은 사건(행 이동·컨틴전시 집행·이슈 개설)
\end{itemize}
\keymsg{에스컬레이션 단계마다 응답 시한이 있다 — 1단계 3영업일 · 2단계 10영업일. IS-05의 8월 CCB 보고 누락이 감리 8번이 됐다}
\note{교재 4.4, 워크북 §3.5 항목 2·3·8(표는 11건 중 6건 발췌). 이슈 로그의 열 규칙 — 발생일과 인지일을 따로(IS-02는 05-12 발생, 05-14 인지), 소유자와 조치자를 따로, 기한과 에스컬레이션 단계를 따로. 리스크 \ra{} 이슈 전환 규칙: 트리거 발동은 신호이므로 소유자가 관측 사실과 일자를 등록부에 적고 집행 준비를 하지만 컨틴전시를 쓰지 않는다. 실현은 확률 100\%가 된 사건이므로 행을 "실현"으로 옮기고 이슈를 개설하며 컨틴전시를 집행한다 — 이때 CCB 보고(RACI 19행 I)가 절차인데 IS-05 R-03 0.50 집행이 8월 CCB에서 빠졌고 그것이 감리 AF-1-08이 되어 10-15 소급 보고한다. 절차 하나가 빠지면 다른 산출물에서 지적이 된다는 예다. 예상 소요 6분.}
\end{frame}

\begin{frame}{리스크 등록부 v1.0 \ra{} v1.2 — 대응은 노출을 없애지 않고 옮긴다}
\fig[0.56]{../그림/PM_Day3/fig_4_4_risk_issue_transitions.pdf}
\figcap{그림 4-4. 종결·실현·이관 4건(R-03·04·09·12), 상승 R-10(0.38 \ra{} 0.76), 신규 7건 R-14\til{}R-20 — 그중 5건(R-14·15·16·17·18)은 대응 조치가 만든 2차 리스크}
\keymsg{R-01 분리 \ra{} R-14 두 번 컷오버, R-04 이연 \ra{} R-17 2027 자금 봉우리 — 관리예비비 대상 2.20은 숫자만 같고 내용이 바뀌었다}
\note{교재 4.5 "2차 리스크의 승격", 워크북 §3.5 항목 4·5. 재평가는 세 숫자를 다시 적는 일이다 — 확률·영향·EMV. R-10(G2 후 마스터 변경)은 9월에 신규 브랜드 상품 코드·B2B 거래처 코드 변경 요구가 접수되어 10\ra{}20\%로 상승했고 컨틴전시를 0.40에서 0.80으로 증액했다. 신규 7건은 출처가 모두 기록에 있다 — R-14는 G1 R-01 분리 의결의 결과(프로그램 PMO장 소유), R-16은 R-03 폴백 배치와 CR-09 22:00 마감의 충돌, R-19는 EX-01 회복 계획(통합시험 압축 \ra{} 결함 유출), R-20은 감리 귀속 분쟁(IS-07). R-17은 원인이 스폰서 자금 결정이므로 컨틴전시가 아니라 관리예비비 대상이다. Day2에서 R-14\til{}R-18을 예비 번호로 남겨 둔 이유가 여기서 드러난다. 예상 소요 5분.}
\end{frame}

\begin{frame}{잔여 EMV 7.26 vs 8.0 — 이내지만 경고, 그리고 컨틴전시 소진 13.8\%}
\fig[0.52]{../그림/PM_Day3/fig_4_5_emv_contingency.pdf}
\figcap{그림 4-5. 6.78 $-$ 종결 1.98 + 신규 3.13 $-$ 재평가 0.67 = 7.26(여유 0.74). 컨틴전시 10.20 = 배정 5.30 + 미배정 2.47 + 집행 1.41 + BL-02 편입 1.02}
\keymsg{여유 0.74는 R-10 하나(0.76)의 크기 — 11월 갱신에서 R-10이 실현되면 즉시 8억 초과. 예외 보고 불요, 스폰서 사전 예고}
\note{교재 4.5 "컨틴전시 소진 추적", 워크북 §3.5 항목 6·7·8. 판정 규칙은 둘 — 잔여 위협 EMV 7.26 < 8.0, 단일 최대 0.80 < 2.0이므로 리스크 축 예외 보고는 불요. 그러나 여유가 1.22에서 0.74로 줄었고 R-10의 트리거(정본 판정 후 변경 요청)는 G2 직후 관측되므로 10-05 성과보고 제9호에 사전 예고한다. 컨틴전시 소진율 13.8\%(약정 1.41/10.20) vs 원가 진척 50.4\%(EV/BAC)는 얼핏 "여유롭다"로 읽히지만, 배정 재조정 후 미배정 풀이 스프린트 4회분에서 1.75회분(2.47)으로 줄었다 — 다음 신규 리스크 2건이면 풀이 마른다. 집행 1.41 중 AC 발생은 1.19(CR-07 0.22는 약정)라는 것, BL-02로 편입된 1.02는 더 이상 컨틴전시가 아니라는 것도 짚는다. 일정 축은 임계경로 잔여 영향 25.5영업일 vs 버퍼 15\til{}30으로 경고다. 예상 소요 6분.}
\end{frame}

\begin{frame}{성과보고 제9호 — RAG 일곱 축과 결정 요청 셋, 그리고 나쁜 소식의 경제학}
\fig[0.44]{../그림/PM_Day3/fig_5_2_mum_effect.pdf}
\figcap{그림 5-2. 일정 R · 원가 A · 범위 G · 품질 A · 리스크 A · 편익 G · 지속가능성 A — 결정 요청: EX-01(10-16) · CCB(10-15) · 감리 귀속 회신(10-08)}
\keymsg{RAG는 색이 아니라 문장이다 — "R"은 "결정이 필요하다"이지 "PM이 못했다"가 아니다. 그 등식이 깨진 조직에서 mum effect가 시작된다}
\note{교재 5.1\til{}5.2, 워크북 §2.5 항목 12(Highlight Report 1장). 일곱 축의 RAG를 학생이 먼저 매기게 하고(4교시 산출물 갤러리워크에서 이미 봤다) 예시본과 비교한다 — 일정 R은 IEAC(t) > 허용이므로 논쟁이 없지만 원가 A/G, 품질 A/R는 조마다 갈린다. 그 갈림이 "판정 기준을 문서에 적었는가"의 문제라는 것이 요점(워크북 §2.3 항목 12 판정 규칙). mum effect(Keil·Smith 계열 연구): 나쁜 소식은 보고자가 책임 귀속을 두려워할수록 늦어지고, 늦어질수록 비싸진다 — 온담에서는 IS-02(설계 승인 대기 12영업일)가 그 사례. 대응은 "보고자의 안전"이 아니라 "판정 기준의 사전 합의"다. 예상 소요 6분.}
\end{frame}

\begin{frame}{동시기록 — 구두 검수 보류 통보에 서면으로 답하지 않으면 발주자 사유 기간을 공제받지 못한다}
\fig[0.44]{../그림/PM_Day3/fig_5_5_contemporaneous_records.pdf}
\figcap{그림 5-5. 구두 통보(09-24) \ra{} 서면 회신(09-26, 발주사 사유 8영업일·R2 인수 FP 12\% 영향 명시) \ra{} 검수 기한 10-16 재확인 \ra{} 회의록 서명}
\keymsg{회신은 논쟁이 아니라 기록이다 — 사실(무엇을, 언제, 누가)과 영향(며칠, 몇 FP)만 적고 책임은 적지 않는다}
\note{교재 5.4 "동시기록 문화", 워크북 §5.5 항목 3(검수 보류 서면 회신 완성본). 세 가지를 짚는다. (1) 구두 통보 자체는 유효하지만 그 내용과 시점의 기록은 회신에서만 남는다 — 회신이 없으면 3개월 뒤 "보류 통보가 있었는가"부터 다툰다. (2) 회신에 반드시 들어갈 다섯 항목: 통보 사실·시점·사유의 귀속(발주사 승인 대기 8영업일)·영향(R2 인수 FP의 12\%)·재확인 요청(검수 기한). (3) 책임을 적지 않는다 — "귀사 귀책으로"가 아니라 "귀사 승인 대기 기간 8영업일"이라고 쓴다. 이것이 SCL 동시지연 분해(다음 프레임)의 재료가 되고, Day1 회의록 서명 회피(권한 상실 신호)와 같은 계열. 6교시 실습 ⑮의 과제. 예상 소요 5분.}
\end{frame}

\begin{frame}{동시지연 분해 — 발주 7 · 동시 3 · 수행 3 · 회복 $-$3, 그리고 지체상금의 계산 순서}
\fig[0.44]{../그림/PM_Day3/fig_5_6_concurrent_delay.pdf}
\figcap{그림 5-6. A10 착수 지연 10영업일의 귀책 분해 — 발주사(설계 승인 대기 IS-02) 7 · 동시(양측) 3 · 수행사(R2 이월) 3 · 수행사 회복 조치 $-$3 = 순 10}
\keymsg{지체상금은 "며칠 늦었나"가 아니라 "며칠 중 며칠이 누구 것인가"에서 시작한다 — 분해가 없으면 계약은 총일수로 청구한다}
\note{교재 5.5 "동시지연", 워크북 §5.5 항목 6. SCL Delay and Disruption Protocol 2판의 동시지연 개념을 온담 숫자로: A10 $-$10 중 발주사 사유(IS-02 설계 승인 대기 12영업일 중 임계경로에 걸린 7)·동시 3(양측 원인 겹침)·수행사 3(R2 이월 420FP)·수행사 회복 $-$3(S18 3주 연장). 지체상금 조항(일 0.05\%, 상한 10\%)은 Day1 계약 구조에서 왔고, 발주사 사유 기간은 공제되므로 순 수행사 귀책은 3 $-$ 3 = 0 — 단 이것은 서면 회신(앞 프레임)이 있을 때만 성립한다. 여유 소유권 문안(Day2 §3, "총여유는 프로젝트의 것")이 여기서 실제로 작동한다. 학생 질문 예상: "동시 3은 누가 정하나" — 답: 계약 문안 + 동시기록, 없으면 감정인. 예상 소요 5분.}
\end{frame}

\begin{frame}{수행사 성과 · 인력 교체 · 하도급 — 계약은 관계가 아니라 기록으로 관리된다}
\fig[0.44]{../그림/PM_Day3/fig_5_4_supplier_performance.pdf}
\figcap{그림 5-4. 스프린트별 velocity(S7\til{}S18)와 결함 유출, SLA 양방향(발주사 승인 SLA 포함), 인력 교체 1건(07-20 인터페이스 PL) 승인, 하도급 재위탁 4 \ra{} 5}
\keymsg{성과 기록이 양방향이어야 협상이 된다 — 수행사 velocity만 적고 발주사 승인 SLA를 안 적으면 IS-02는 수행사 탓이 된다}
\note{교재 5.3, 워크북 §5.5 항목 2·4·5. 세 기록의 요점. (1) 수행사 성과: velocity와 결함 유출률을 스프린트마다 기록하되, 발주사 측 SLA(승인 대기일·자료 제공일)를 같은 표에 둔다 — 그래야 S10\til{}S11 하락의 원인이 IS-02(발주사)임이 표에서 읽힌다. (2) 인력 교체: 07-20 인터페이스 PL 교체 승인서 — 교체 사유·인수인계 기간·품질 영향 평가가 없으면 나중에 R-15(선행 릴리스 결함)의 귀책을 못 가른다. (3) 하도급: 재위탁 4 \ra{} 5 승인 기록 — 개인정보 위탁·재위탁(최영주 CPO)과 연결되어 Day4 개인정보 처리 위탁 검토서의 입력이 된다. 계약 변경 로그 10건은 ⑮ 완성본. 예상 소요 5분.}
\end{frame}

\begin{frame}{5교시 핵심 정리}
\begin{enumerate}
\item 결함밀도 0.132 > 0.13 — 분모(인수 FP)와 분자(시험 결함)의 정의가 먼저, 판정은 그다음
\item 8월 평탄은 개선이 아니라 시험 정지(TC 210 \ra{} 95) — 곡선은 반드시 시험 노력과 함께 읽는다
\item 감리 지적은 등급보다 원인 귀속 칸 — 3·4번 이의 서면이 나중의 책임 귀속을 정한다
\item 이슈 11건·리스크 v1.2·잔여 EMV 7.26 < 8.0(여유 0.74 = R-10 하나) — 이내지만 경고, 사전 예고
\item 서면 회신·동시지연 분해·양방향 성과 기록 — 계약은 기록으로 관리된다
\end{enumerate}
\keymsg{6교시 실습 — 이슈 2건·리스크 재평가·감리 귀속 이의·검수 보류 회신·변경 대가를 직접 쓴다}
\note{교재 제4장·제5장 핵심 정리. 생각해 볼 질문(제5장 2 — "여러분 회사의 마지막 검수 보류는 서면으로 남아 있는가"). 6교시 준비물(워크북 §3.4·§4.4·§5.4 빈 템플릿, 감리 지적 원문 8건 사본)을 안내. 예상 소요 3분.}
\end{frame}

% ===== 6교시 (75분) — 실습 ② 이슈 로그·품질 기록·감리 대응·검수 보류 회신 =====
\section{6교시 (75분) — 실습 ② 이슈 로그·품질 기록·감리 대응·검수 보류 회신}

% =====================================================================
% 6교시 — 실습 ②
% =====================================================================
\begin{frame}{6교시 — 이 교시에서 하는 것}
\begin{itemize}
\item 5교시의 세 문서(⑬·⑭·⑮)를 \textbf{손으로} 쓴다 — 워크북 §3\til{}5
\item 세 과제, 한 사건 — 감리 지적 3·4번(IS-07)이 검수 보류·대기 비용·R-20으로 번진다
\item ⑬ 이슈 2건 신규 기입 + 리스크 1건 재평가 (25분)
\item ⑭ 감리 지적 3·4번 원인 귀속 이의 서면 (20분)
\item ⑮ 구두 검수 보류에 서면 회신 + CR-08·11 변경 대가 (25분)
\end{itemize}
\keymsg{세 문서는 같은 사건을 세 자리에서 기록한다 — 귀속 칸 하나가 세 문서의 숫자를 정한다}
\note{4교시 실습 ①과 같은 운영 — 개인 작성 후 조별 대조, 강사는 순회. 다른 점은 세 과제가 한 사건(09-25 감리 1차 지적 8건, 그중 3·4번의 "수행사 설계 미흡" 기재)에서 갈라진다는 것이다. ⑬은 이 사건이 이슈 로그(IS-07)와 등록부(R-20)에 어떻게 들어가는지, ⑭는 귀속을 어떻게 뒤집는지, ⑮는 그 결과가 검수와 대기 비용에 어떻게 이어지는지를 쓴다. 시간이 부족하면 ⑮의 변경 대가 산정표는 두 행(CR-08·CR-11)만 쓰게 한다. 워크북 §3.7·§4.7·§5.7 기본 과제와 다르게, 오늘 교실에서는 완성 예시본이 다루는 사건 자체를 재구성한다(예시본 숫자를 베끼지 않고 근거 문서에서 다시 도출). 예상 소요 3분.}
\end{frame}

\begin{frame}{과제 ⑬ — 워크북 §3: 이슈 2건 신규 기입 + 리스크 1건 재평가}
\begin{itemize}
\item 시점 2026-11-30, G2 직후 — 워크북 §3.7 기본 과제 상황
\item IS-12: 영업본부 "델리 서브 브랜드 상품 코드 체계 변경" 정식 요구 (R-10 실현)
\item IS-13: P2 벤더 MC 2027-03-19로 2주 추가 지연 통보 (R-01 잔여 실현)
\item 재평가: R-10 행 — 확률·영향·EMV·컨틴전시 0.80과 부족분 자금원
\item 항목 6 잔여 EMV 대조 — 7.26 \ra{} ? / 허용범위 8.0 넘는가
\end{itemize}
\keymsg{이슈 행은 열 아홉 개를 다 채운다 — 특히 "출처"(R-nn 전환 / 신규 발견)와 "기한"·"에스컬레이션 단계"}
\note{워크북 §3.3 항목 2(이슈 로그 열 규칙)와 항목 3(리스크 \ra{} 이슈 전환 규칙)을 펴게 한다. IS-12는 R-10 실현이므로 등록부 R-10 행을 "실현 \ra{} IS-12"로 바꾸고, 재발 가능성이 있으면 같은 번호로 잔여 리스크를 유지한다(R-05 선례, §3.5 항목 4). R-10의 3.8억 전액 실현은 컨틴전시 배정 0.80으로 부족하므로 부족분 자금원 순서(미배정 풀 \ra{} 관리예비비)와 CFO 병행 결재 조건(5억 초과 아님, 5,000만 원 이상 내부통제)을 적게 한다. IS-13은 R-01 잔여 실현이자 관리예비비 대상이므로 당일 스폰서·프로그램 PMO장 보고(Day2 §5.5 항목 9)가 에스컬레이션 칸에 들어가야 한다. 잔여 위협 합이 8.0을 넘으면 예외 보고 대안 3개(대응 강화 / 관리예비비 이관 / 허용범위 재협상)와 결정자(RACI 20·21행)까지. 예상 소요 4분(설명) + 25분(작성).}
\end{frame}

\begin{frame}{과제 ⑭ — 워크북 §4: 감리 지적 3·4번 원인 귀속 이의 서면}
\begin{itemize}
\item AF-1-03 (중대): P2 설비 인터페이스 규격 설비측 확인 서명 부재 — 감리 기재 "수행사"
\item AF-1-04 (보통): 가맹 SW 배포 범위(1.1.6) 설계 미확정 — 감리 기재 "수행사"
\item 대응 관리대장 항목 6의 두 행을 채운다 — 발주사 판정 귀속·근거·조치·기한·이의
\item 이의 서면 1장 — 수신 수석감리원, 참조 스폰서·PMO장·수행사 PM
\item 근거는 두 곳뿐 — WBS 담당 열(발주/수행/공동)과 RACI 행
\end{itemize}
\keymsg{지적 사실은 인정하고 귀속만 다툰다 — 등급·내용에 대한 이의가 아님을 첫 문장에 쓴다}
\note{워크북 §4.3 항목 6과 §4.5 항목 6의 대장 양식을 쓴다. AF-1-03의 근거는 규격이 WBS 1.3.1 발주 담당 산출물이고 P2 전달 책임이 RACI 9행 프로그램 PMO장(A)이라는 것, 그리고 설비 벤더의 05-29 변경 요청에 대한 PMO 회신이 미결이라는 사실 기록(회신 요청 06-05·07-10·08-21, 수행사 서면 통지 06-08·08-24). AF-1-04는 1.1.6 대응안 결정이 RACI 11행 오정근(A)이며 05-08 회신으로 결정을 유보했다는 기록. 서면의 네 단락 — 사실 인정 / 귀속 재분류 요청과 근거 / 발주사 조치 계획(10-16 규격 회신, 11-27 대응안 결정) / 요청 사항(정본 정정, 10-08 회신, 향후 귀속 판정 기준 합의). 교재 §4.3 "왜 귀속 두 글자가 등급보다 무거운가"를 상기시킨다. 예상 소요 3분(설명) + 20분(작성).}
\end{frame}

\begin{frame}{과제 ⑮ — 워크북 §5: 구두 검수 보류에 서면 회신 + 변경 대가}
\begin{itemize}
\item 사건 — 09-28 주간회의 구두 "감리 지적 정리 후 검수" \ra{} 수행사 09-29 서면 "발주사 사유 보류, 대기 비용"
\item 회신 서면 1장(항목 4) — 검수는 §14대로 진행 중, 보류가 아님을 먼저
\item 발주사 사유 기간을 발주사가 스스로 적는다 — 09-28\til{}10-08, 8영업일, 영향 범위 R2 FP 12\%
\item 변경 대가 산정표(항목 2) — CR-08 +90 FP / CR-11 $-$60 FP, 단가 55만 원
\item 항목 6 귀책 분해표에 한 행 — "AF-1-03·04 귀속 정리 8일, 발주사 단독, 병행 가능"
\end{itemize}
\keymsg{구두에는 서면으로, 보류에는 조건으로, 대기 비용에는 기간·범위·판정 절차로 답한다}
\note{워크북 §5.3 항목 2·4·6과 §5.5 완성본을 쓴다. 회신의 뼈대 — ① 검수 요청 접수 사실과 §14 기산일(09-25)·통보 기한(10-16) ② 구두 통보의 정정(보류가 아님) ③ 검수 완료 조건(수행사 귀속 3건 조치) ④ 발주사 사유 기간의 명시(8영업일, 관련 산출물 한정 약 12\%, 나머지와 R3 병행 가능) ⑤ 대기 기록 10-09 제출 \ra{} 10-14 변경협의체 판정 ⑥ 요청 사항. 변경 대가는 CR-08 +90 × 55만 = +0.495억(R4 검수 후 지급), CR-11 $-$60 × 55만 = $-$0.33억(감액, 잔금 차감) — 감액도 합의서 6호에 같은 행으로. 시간이 남는 조는 CR-11의 데이터 전환 용역 $-$2.00(DT-01)까지 항목 7 계약 변경 로그에 넣게 한다. 예상 소요 4분(설명) + 25분(작성).}
\end{frame}

\begin{frame}{작성 요령 ① — 원인 귀속 칸이 책임 귀속을 정한다}
\fig[0.58]{../그림/PM_Day3/fig_4_3_audit_attribution.pdf}
\figcap{그림 4-3. 감리 지적 8건 — 감리 기재 귀속(수행사 6·발주사 2) vs 발주사 판정 귀속(수행사 3·발주사 4·공동 1). AF-1-03·04의 이의 근거는 WBS 담당 열과 RACI 행 두 가지뿐이다}
\keymsg{감리 기재를 그대로 받은 대장은 검수·변경협의체에서 발주사의 미결정을 수행사의 미흡으로 기록한다}
\note{교재 §4.3. 대장에 "감리 기재 귀속"과 "발주사 판정 귀속" 두 열을 따로 두는 이유를 그림으로 보인다 — 한 열이면 감리 기재가 곧 판정이 된다. 판정 근거 열에는 WBS 사전 카드의 담당(발주/수행/공동)과 RACI 행 번호, 그리고 일자가 붙은 사실 기록(발송·회신 요청·서면 통지)만 적는다. 형용사("명백히", "당연히")는 근거가 아니다. 8건 중 이의는 2건뿐이고 나머지 6건은 그대로 받는다 — 전부 다투면 신뢰를 잃는다. 이 두 열의 차이가 ⑮의 귀책 분해표(발주사 단독 8일)와 R-20의 EMV로 이어진다는 것을 짚는다. 예상 소요 4분.}
\end{frame}

\begin{frame}{작성 요령 ② — 서면 회신에 반드시 들어갈 것}
\begin{tbl}[\scriptsize]{L{3.2cm}L{5.8cm}L{2.6cm}}
항목 & 온담 R2 회신(10-01)에서 & 없으면 \\
\midrule
접수 사실·계약 조항·기한 & 09-25 접수, §14 15영업일, 통보 기한 10-16 & 기산일 분쟁 \\
구두 통보의 정정 & "보류"가 아니라 "진행 중, 조건부" & 수행사 기록만 남음 \\
완료 조건 & 수행사 귀속 3건 조치(10-16\til{}10-30) & 검수 무기 연기 \\
발주사 사유 기간 & 09-28\til{}10-08 8영업일, 범위 R2 FP 12\% & 전 기간·전 인력 청구 \\
판정 절차·시한 & 대기 기록 10-09 \ra{} 변경협의체 10-14 & 협상 근거 없음 \\
\end{tbl}
\keymsg{발주사가 스스로 "8영업일"을 적는 이유 — 기간과 범위를 먼저 정의하는 쪽이 협상의 자를 쥔다}
\note{교재 §5.5 "회신의 구조"와 "왜 발주사가 스스로 8영업일을 적는가". 수강생이 가장 저항하는 부분이 넷째 행이다 — 발주사 사유를 왜 발주사가 먼저 인정하느냐. 답은 그림 5-5의 기록 A(수행사만 기록)와 기록 B(발주사 회신)의 차이다. 기록 A만 있으면 대기 기간은 수행사가 정한 기간 전체(09-28부터 무기한)와 전 인력이 되고, 기록 B가 있으면 8영업일 × 관련 산출물 12\%로 한정된다. 인정하지 않는 것이 아니라 한정하는 것이다. 참조 수신자(구매팀장·재경팀장·수석감리원)를 넣는 이유도 같다 — 회신이 계약·지급·감리 세 기록에 동시에 들어간다. 예상 소요 4분.}
\end{frame}

\begin{frame}{작성 요령 ③ — 변경 대가 = FP 델타 × 55만 원, 감액도 같은 표에}
\begin{tbl}[\scriptsize]{L{3.0cm}rrL{1.6cm}L{2.2cm}L{2.0cm}}
CR / 건 & FP $\Delta$ & 금액(억) & 증/감 & 합의·서면 & 지급 시점 \\
\midrule
CR-08 파일럿 임시 연동 & +90 & +0.495 & 증액 & 09-30 · 6호 & R4 검수 후 \\
CR-11 5년 초과 이력 제외 & $-$60 & $-$0.33 & \textbf{감액} & 09-30 · 6호 & 잔금 차감 \\
데이터 전환 용역(CR-11) & — & $-$2.00 & 감액 & 09-25 · DT-01 & 11.4 \ra{} 9.4 \\
\end{tbl}
\begin{itemize}
\item 네 가지가 다 있어야 한 행 — FP 산정 근거 · 단가 · 변경협의체 합의 일자 · 서면 번호
\item FP 확인 3인 서명(PMO 기준선 담당·수행사 PM·형상 담당) — CR-07\til{}11은 BL-02 확인서에서 일괄
\item 대기 비용은 FP가 아니라 인일 × 부속서 단가 — 별도 열, 같은 표
\end{itemize}
\keymsg{감액을 변경으로 안 쓰면 G4 정산에서 수행사가 "합의한 적 없다"고 말할 수 있다}
\note{교재 §3.5와 워크북 §5.3 항목 2. 고정가 FP 계약에서 변경 대가는 산식이 하나뿐이라 쉬워 보이지만, 넷 중 하나가 빠진 행이 G4에서 다시 협상된다(§5.7 점검 질문 1). CR-11은 SI 계약에서 $-$60 FP 감액이면서 데이터 전환 용역에서 $-$2.00 감액이라 두 계약·두 서면(SI 합의서 6호·DT-01)에 각각 들어간다 — 한 변경이 두 행이다. 대기 비용(R-05 0.42, 이번 AF-1-03·04 건은 10-14 판정 후)은 FP 델타가 아니므로 "대기 비용" 열을 따로 두되 같은 표에 쓴다 — 그래야 누적(SI +1.4375 = FP +185 · 대기 0.42)이 한 줄로 나온다. 예상 소요 3분.}
\end{frame}

\begin{frame}{흔한 오류 — 실습 ②에서 반복되는 네 가지}
\begin{itemize}
\item \textbf{이슈를 리스크 등록부에 넣는다} — 실현된 것은 확률 100\%가 아니라 이슈, 등록부 행은 "실현 \ra{} IS-nn"
\item \textbf{귀속 칸을 비운다} — "추후 협의"는 감리 기재가 판정이 되는 것과 같다
\item \textbf{구두 통보에 구두로 답한다} — 수행사 서면(09-29)만 남고 발주사 기록은 0건
\item \textbf{감액을 변경으로 안 쓴다} — CR-11 $-$60 FP가 로그에 없으면 잔금 차감 근거가 없다
\item 기한 없는 이슈 — 이슈가 아니라 배경이다 (워크북 §3.7 점검 질문 3)
\end{itemize}
\keymsg{네 오류의 공통점 — 기록해야 할 순간에 기록하지 않아 상대방의 기록이 사실이 된다}
\note{순회 중 발견한 오류를 이 프레임에서 모아 말한다. 첫째는 §4.4 전환 규칙 — 트리거 발동은 신호(등록부에 관측 사실 기록), 실현은 사건(이슈 로그 등재 + 컨틴전시 집행 개시). 둘째는 귀속 칸을 "공동"으로 얼버무리는 변형도 포함한다 — AF-1-01처럼 근거가 있어 공동이면 되지만 근거 없는 공동은 빈칸이다. 셋째는 mum effect(§5.2)의 계약판 — 나쁜 기록을 남기기 싫어 서면을 피하면 상대의 서면이 유일한 기록이 된다. 넷째는 CR-11이 "좋은 소식"이라 로그를 건너뛰는 경우 — 감액도 변경이고(§3.5), 서면 없는 감액은 G4에서 사라진다. 다섯째는 덤 — 기한과 다음 에스컬레이션 단계가 없는 이슈 행은 완성이 아니다. 예상 소요 4분.}
\end{frame}

\begin{frame}{예시본 참조 — 워크북 완성 예시본과 템플릿 PDF}
\begin{tbl}[\scriptsize]{L{1.2cm}L{3.4cm}L{3.2cm}L{3.6cm}}
산출물 & 워크북 & 빈 템플릿 PDF & 완성 예시 PDF \\
\midrule
⑬ & §3.5 (이슈 11건·재평가·EMV 7.26) & 13\_이슈리스크갱신\_빈템플릿 & 13\_이슈리스크갱신\_완성예시 \\
⑭ & §4.5 (대장 8건·이의 서면 09-29) & 14\_품질통제기록\_빈템플릿 & 14\_품질통제기록\_완성예시 \\
⑮ & §5.5 (산정표·회신 10-01·귀책표) & 15\_조달계약변경\_빈템플릿 & 15\_조달계약변경\_완성예시 \\
통합 & — & 00\_Day3\_템플릿팩\_빈양식 (46쪽) & 00\_Day3\_템플릿팩\_완성예시 (42쪽) \\
\end{tbl}
\begin{itemize}
\item 경로 — \textsf{교재/템플릿/PM\_Day3/}, 교재·워크북과 다르면 \textbf{워크북 완성 예시본이 정본}
\item 예시본은 대조용이다 — 작성 후에 편다. 먼저 펴면 베끼게 된다
\end{itemize}
\keymsg{예시본과 내 답이 다른 자리에 배움이 있다 — 다른 이유를 한 줄로 적어 두라}
\note{작성 시간이 끝난 뒤(약 70분 시점) 예시본을 펴게 한다. 대조 순서 — ⑬은 §3.5 항목 2의 IS-07 행과 항목 5의 R-20 행(오늘 사건이 어떻게 두 문서에 들어갔는가), 항목 6의 변동 분해(6.78 $-$ 1.98 + 3.13 $-$ 0.67 = 7.26)로 자기 재평가의 산식을 검산. ⑭는 §4.5 항목 6의 AF-1-03·04 두 행과 09-29 서면의 네 단락 — 자기 서면에 "지적 사실의 인정"이 첫 문장에 있는가. ⑮는 §5.5 항목 4 회신의 4항(발주사 사유 기간)과 항목 6 귀책 분해표의 "발주사 단독 8일, 병행 가능, 임계경로 영향 0". 가로 페이지 표(이슈 로그·리스크 재평가·감리 대응 대장·변경 대가·귀책 분해)는 템플릿팩 인쇄본을 쓴다. 예상 소요 5분(대조).}
\end{frame}

\begin{frame}{시간 배분과 심화 과제}
\begin{tbl}[\scriptsize]{L{1.4cm}L{6.2cm}L{2.6cm}}
분 & 활동 & 산출 \\
\midrule
0\til{}5 & 과제 설명 (프레임 2\til{}4) & — \\
5\til{}30 & ⑬ IS-12·IS-13 기입, R-10 재평가, EMV 대조 & §3 항목 2·4·6·7 \\
30\til{}50 & ⑭ 대장 2행 + 이의 서면 1장 & §4 항목 6 \\
50\til{}70 & ⑮ 회신 서면 1장 + 변경 대가 2행 + 귀책 1행 & §5 항목 2·4·6 \\
70\til{}75 & 예시본 대조·정리 & 다른 자리 한 줄 \\
\end{tbl}
\begin{itemize}
\item \textbf{심화 과제(3\til{}7년차)} — 자기 회사 이슈 20건의 출처(등록부에 있었나)와 소유자가 PM인 비율을 세라 — 워크북 §3.7
\item 계약서에서 검수·지체상금·변경 세 조항의 위치를 한 장으로 — 워크북 §5.7
\end{itemize}
\keymsg{75분에 세 문서 여섯 조각 — 완성보다 빈칸이 어디인지 아는 것이 오늘의 산출이다}
\note{시간이 밀리면 ⑮의 변경 대가 두 행을 숙제로 돌리고 회신 서면을 지킨다 — 서면 회신이 이 교시의 핵심 산출이기 때문이다. 조별 대조는 4교시처럼 옆 조와 교환해 ⑭ 서면을 상대 조가 "수행사 PM" 역할로 읽고 반박하게 하면 귀속 근거의 약한 자리가 드러난다(3분). 심화 과제 둘은 워크북 §3.7·§5.7의 것을 그대로 — 등록부에 있었던 이슈 비율 30\% 미만이면 등록부가 위험을 못 보는 것, 90\% 이상이면 이슈 로그가 등록부 복사본. 계약 세 조항 요약은 Day4 종료·정산과 연결된다. 예상 소요 2분.}
\end{frame}

\begin{frame}{6교시 핵심 정리}
\begin{itemize}
\item 실현된 리스크는 이슈다 — 등록부 행은 "실현 \ra{} IS-nn", 재발 가능성만 같은 번호로 남긴다
\item 잔여 EMV는 분해해서 본다 — 종결 $-$ / 신규 + / 재평가 ± (6.78 \ra{} 7.26, 여유 0.74)
\item 감리 대장의 "발주사 판정 귀속" 열 — 근거는 WBS 담당 열과 RACI 행, 사실은 일자와 함께
\item 구두 검수 보류에는 서면으로 — 조항·기한·조건·발주사 사유 기간·판정 절차 다섯 가지
\item 변경 대가 한 행 = FP $\Delta$ × 55만 + 합의 일자 + 서면 번호, 감액도 같은 표
\end{itemize}
\keymsg{통제의 마지막 층은 기록이다 — 누가 언제 무엇을 결정하지 않았는지를 발주사가 먼저 적는다}
\note{Day3 산출물 5종이 모두 결정 요청으로 끝난다는 도입의 한 문장으로 돌아간다. ⑬의 결정 요청은 컨틴전시 재배정과 관리예비비 이관(R-10 부족분·R-14·R-17), ⑭는 귀속 정정 회신 10-08과 발주사 조치 2건(10-16·11-27), ⑮는 변경협의체 10-14의 대기 비용 판정과 합의서 6호 서명. 셋 다 결정자가 발주사 계층(프로그램 PMO장·오정근·스폰서)이라는 점 — 수행사에 요구하기 전에 발주사 안의 미결정을 기록하는 것이 Day3 오후의 주제였다. 마무리 30분(제6장 학술 배경·읽을 자료·Day4 예고)으로 넘어간다. 예상 소요 3분.}
\end{frame}

% ===== 마무리 (30분) — 학술 배경 · 읽을 자료 · Day4 예고 =====
\section{마무리 (30분) — 학술 배경 · 읽을 자료 · Day4 예고}

% =====================================================================
% 마무리 (30분) — 학술 배경 · 읽을 자료 · Day4 예고 — 교재 제6장
% =====================================================================
\begin{frame}{마무리 — 이 교시에서 배우는 것}
\begin{itemize}
\item 오늘 인용한 문헌 아홉 갈래 중 여섯 — "왜 그런가"와 "언제 틀리는가"
\item 읽을 자료 — 필수 4종(과정 전)·심화 13종(30일·90일)·참고 9종
\item 오늘 만든 산출물 ⑪\til{}⑮ 다섯 — 서로를 \textbf{설명하는} 구조 한 장
\item Day2 문서로 되돌아간 것 12건 — 재기준선은 0회
\item Day4 예고 — 전환·종료 ⑯\til{}⑳, 그리고 사후 읽기 과제
\end{itemize}
\keymsg{표준은 "이렇게 하라"까지만 말한다 — 지표는 왜 정보를 잃고, 나쁜 소식은 왜 올라오지 않는가는 문헌이 말한다}
\note{교재 제6장 도입. Day1 제4장·Day2 제6장과 같은 자리다 — 실행·통제 단계에서 학술 문헌이 답하는 두 질문은 지표(왜 SPI는 종료에서 1이 되는가)와 사람(왜 나쁜 소식은 올라오지 않는가)이다. 30분 배분 — 학술 배경 12분, 읽을 자료 5분, 산출물 5종 되짚기·Day2 갱신 6분, Day4 예고·과제 7분. 6교시 실습 ②가 늦게 끝났으면 학술 배경 세 프레임을 각 2분으로 줄이고 Day4 예고는 자르지 않는다 — Day4 첫 시간의 입력이 오늘 문서이기 때문이다. 서지의 [검증 필요] 표기는 과정 리딩리스트(05\_리딩리스트)의 것을 그대로 따랐고, 슬라이드에는 표기를 생략했으니 인쇄 전 원문 대조를 권한다고 말해 둔다. 예상 소요 2분.}
\end{frame}

\begin{frame}{학술 배경 ① — 지표는 왜 정보를 잃는가: Lipke · Christensen}
\begin{tbl}[\scriptsize]{L{2.3cm}L{4.2cm}L{4.6cm}}
문헌 & 발견 & 온담에서 \\ \midrule
Lipke 2003 \textit{Schedule is Different} & SV·SPI는 원가 단위 \ra{} 종료 시 반드시 SV = 0·SPI = 1, 후반 일정 지표는 정보 상실. 해법 = EV를 PV 곡선에 투영해 시간축 ES & ES 8.28개월 · SV(t) $-$0.72 · SPI(t) 0.920 vs SPI 0.943 — 벌어지기 시작한 7월(진척 40\%)이 가장 정보가 많다 \\
Christensen \& Payne 1992 · Christensen \& Heise 1993 & 미 국방 원가정산 계약 표본 — 누적 CPI는 진척 20\% 이후 ±10\% 안정. "CPI는 좋아지지 않는다" & 표본이 다르다(고정가·발주사 관점). 쓸모는 희망 꺾기 — 조정 CPI 0.968이 50\%에서 안정이면 PMB 초과 확정, 문제는 집행 한도 156.0 안인가 \\
\end{tbl}
\keymsg{흔한 오용 둘 — ES를 CPM의 대체로 읽는 것(지연의 위치는 CPM만 안다) / 착시 포함 CPI 1.032를 20\% 규칙으로 외삽하는 것}
\note{§6.1·§6.3. Lipke의 원 논문은 네 쪽짜리 실무 논문(The Measurable News 2003)이고 EVM 40년의 결함 하나를 고쳤다 — 후속으로 Lipke의 Earned Schedule(2009)과 Vanhoucke의 실증이 ES 기반 예측이 SPI 기반보다 일관되게 정확함을 보였다. 세 오독을 짚는다 — ES는 EVM의 대체가 아니라 확장(같은 PV·EV) / ES는 프로젝트 평균이라 지연의 위치를 모른다(온담 $-$15영업일 평균 vs 임계경로 $-$10) / 후반부에만 쓰는 것이 아니라 두 지수가 벌어지는 시점이 가장 정보가 많다. Christensen의 20\% 규칙은 EAC1의 실무적 정당화가 되었지만 표본은 원가정산 계약이며, 온담 같은 고정가·발주사 관점 EVM에 그대로 옮기면 틀린다(§2.5). 규칙의 진짜 쓸모는 "CPI가 곧 회복된다"는 희망을 꺾는 데 있고, 가장 흔한 오용은 착시 포함 CPI(1.032)를 외삽하는 것이다. 예상 소요 4분.}
\end{frame}

\begin{frame}{학술 배경 ② — 완료된 것만 재는 지표의 병리: Vacanti · Reinertsen}
\begin{itemize}
\item Reinertsen 2009 — 이용률 \ra{} 100\%이면 대기 시간 폭증(M/M/1의 $\rho/(1-\rho)$와 Kingman 근사는 다른 식)
\item R-05의 구조 — 발주사 14명 겸임 \til{}100\% \ra{} 결정 요청이 큐에 쌓임. 처방은 "더 열심히"가 아니라 이용률 낮추기
\item Vacanti 2015 — WIP·cycle time·throughput·work item age. Little's Law는 예측이 아니라 \textbf{안정성 진단} 도구
\item flow debt — 오래된 항목 방치 + 새 항목 빨리 끝내기 = 평균 개선, 부채 누적. SPI 수렴과 같은 병리
\item 백분위수(50/85/95\%)로 약속한다 — 평균이 아니라 분포. S10 WIP 41·age 12가 SI EV 하락의 선행 신호
\end{itemize}
\keymsg{처방은 같다 — 완료된 것 옆에 진행 중인 것의 나이를 병기한다(§2.7 흐름 지표, §4.2 결함 나이)}
\note{§6.4. Day2 §6.4가 Reinertsen을 배치 크기·큐·WIP의 경제학으로 소개했다면 Day3에서 그것은 측정이 된다. 큐 논증을 인용할 때 M/M/1과 Kingman 근사를 섞지 않도록 주의시킨다. Vacanti의 세 기여 — Little's Law의 재정의(성립 가정이 깨지는 구간이 불안정의 신호 — 원문 대조 필요), flow debt, Cycle Time Scatterplot과 백분위수. 2교시 §2.7 표(S7\til{}S18)는 이 넷을 EVM 옆에 둔 것이고, 5교시 §4.2가 같은 처방(work item age 병기)을 결함 나이에 적용했다. 학생에게 묻기 — "여러분 팀의 번다운은 완료만 재는가, 진행 중 항목의 나이도 재는가". 예상 소요 3분.}
\end{frame}

\begin{frame}{학술 배경 ③ — 사람: Keil·Smith · Staw · Brooks}
\begin{tbl}[\scriptsize]{L{2.6cm}L{4.1cm}L{4.4cm}}
문헌 & 발견 & 이 교재가 빌린 처방 \\ \midrule
Smith \& Keil 2003 · Keil, Im \& Mähring 2007 (mum effect, 원전 Rosen \& Tesser 1970) & 정보 비대칭 · 비난 예상(blame) · 조직 분위기(climate)가 나쁜 소식의 보고를 가늘게 만든다 & 성격 교정이 아니라 구조 — 트리거를 관측 항목으로, 예외 보고자의 대우로 climate를 정한다(§5.2, R-05 05-12\ra{}07-15) \\
Staw 1976 \textit{Knee-deep in the big muddy} & escalation의 핵심은 sunk cost가 아니라 self-justification — 자기가 내린 결정일수록 추가 투입 & 착수 결정자 $\neq$ 중단 결정자 · 대안은 각각 다른 사람의 결정 · 사전에 정한 숫자(11-13 S22 속도 620)가 자동 발동 \\
Brooks 1975 \textit{Mythical Man-Month} & 늦은 프로젝트에 인력을 더하면 더 늦어진다 — 학습·교육 부담 + 경로 $n(n-1)/2$(82명 = 3,321) & 회복 대안 1은 인력 추가가 아닌 병행·재배치 / 교체는 전임·후임 동시 투입 / 결함 전담 2명은 R3 병행 해제로 \\
\end{tbl}
\keymsg{셋 다 결론은 "사람을 고치지 말고 결정의 구조를 미리 정하라"다 — EX-01의 대안 3개가 그 구조다}
\note{§6.5\til{}6.7. Keil·Smith 계열은 소프트웨어 프로젝트에 mum effect를 이식한 일련의 연구이고, 같은 갈래의 Kappelman, McKeeman \& Zhang(2006) Dominant Dozen이 조기 신호 12개를 사람·프로세스 요인으로 순위화해 상위 신호가 대부분 기술이 아님을 보였다 — §2.8 조기 신호 5개의 이론적 틀. Staw의 실험은 Day1이 게이트의 중단 기준을 착수 결정과 동시에 제출하게 하고 판정자를 분리한 근거이며, Day3에서는 예외 보고의 대안이 "각각 다른 사람의 결정"이어야 하는 이유(§4.6)다 — Flyvbjerg 열 편향의 ⑩. Brooks는 Day2 §3.4 crashing의 한계로 이미 인용했고 오늘 세 자리(회복 계획·인력 교체·결함 전담)에 다시 나타났다 — Williams(2002)의 시스템 다이내믹스가 이 격언을 압축\ra{}병행\ra{}재작업\ra{}추가 압축의 정량 구조로 승격시켰다. Kerzner(§6.8)는 참조서 — 90\% 증후군과 스폰서십 실패 유형만 빌린다. 예상 소요 4분.}
\end{frame}

\begin{frame}[shrink=20]{읽을 자료 — 필수(과정 전) · 심화(30일 · 90일)}
\begin{tbl}[\scriptsize]{L{1.2cm}L{4.6cm}L{5.4cm}}
구분 & 자료 & 어느 부분 · 왜 \\ \midrule
필수 & Lipke, \textit{Schedule is Different} (Measurable News, 2003) & 네 쪽 전체 · ES 네 공식을 강의장에서 처음 보면 오후가 무너진다 \\
필수 & ANSI/PMI 19-006-2019 \textit{Standard for EVM} 용어 절 · PRINCE2 7 Progress practice & PMB·BAC·관리예비비 정의 · 허용범위·예외 보고 원문 \\
30일 & Fleming \& Koppelman \textit{EVPM} 4판(PMI 2010) · Christensen 두 논문 · Vacanti 2015 & 고정가 EVM·EAC 전제 · CPI 안정성 표본의 한계 · flow debt·Little's Law \\
30일 & Staw 1976 (\textit{OBHP} 16(1)) · Smith \& Keil 2003 (\textit{ISJ} 13(1)) · Kappelman et al. 2006 & self-justification · 나쁜 소식의 구조 · 조기 신호 12개 \\
90일 & Reinertsen 2009 · SCL Protocol 2판(2017) · PMBOK 8판 M\&C · PRINCE2 7 세 프로세스 · Brooks 『맨먼스 미신』 · Kan 2판 · \textit{Accelerate} + DORA 2025 & 큐·이용률 · 원칙 8·10·14 · 통제 활동 배치 · 발주·수행 PM 분리 · Rayleigh·DRE · 4지표 정의 \\
\end{tbl}
\keymsg{읽지 않아도 되는 것 — 전제 없는 EVM 공식 카드, SPI·CPI 팀 KPI 대시보드 백서, DORA 서열화 사례집, Lewin 3단계}
\note{§6.10. 필수 4종 중 첫째는 사실 Day2 워크북 완성 예시본 5종(60분)이다 — 슬라이드에는 문헌만 실었다. 참고 9종(워크북을 쓰다가 손을 뻗을 것 — ANSI/PMI 19-006 Baseline maintenance, PMI 형상관리 실무표준 2007, PMBOK 6판 7.4.2, PRINCE2 7 Issues·Change practice, Hillson \& Simon ATOM 3판, ISO 21502 §7, 감리기준 고시 2024-53호, SW진흥법·하도급법 조문, SCL 원칙 10)은 교재 표를 가리킨다. "읽지 않아도 되는 것"은 읽은 것이 잘못이 아니라 인용할 때 한계를 함께 말하라는 뜻 — Lewin의 unfreeze–change–refreeze는 Cummings, Bridgman \& Brown(2016)이 Lewin이 만든 적 없음을 사료로 보였다. 감리기준(§6.9)은 민간 온담이 빌린 구조이며 빌릴 때 귀속 판정 원칙을 함께 빌려야 한다는 점만 한 줄 덧붙인다. 예상 소요 4분.}
\end{frame}

\begin{frame}[shrink=12]{오늘 만든 다섯 — 성과보고가 중심이고 나머지 넷이 각주다}
\begin{tbl}[\scriptsize]{L{0.5cm}L{3.2cm}L{3.6cm}L{3.6cm}}
\# & 산출물 (워크북 §) & 오늘의 숫자 & ⑫와의 연결 \\ \midrule
⑪ & 변경요청서·영향 분석·CCB 의결서 (§1) & CR-08\til{}11 · +185 FP · +1.02 · CCB 10-15 · BL-02 146.82 & EAC4 ETC 행 "변경 대가 1.02"의 출처 \\
⑫ & 성과보고 EVM·ES 제9호 (§2) & EV 73.43 / AC 71.13 · SPI 0.943 · SPI(t) 0.920 · CPI 1.032\ra{}0.968 · EAC4 150.86 · A10 $-$10 \ra{} EX-01 & 중심 \\
⑬ & 이슈 로그·리스크 등록부 갱신 (§3) & 이슈 11건 · 잔여 EMV 7.26 · 컨틴전시 10.20 = 5.30+2.47+1.41+1.02 & SV 음수의 이유 — IS-02(R-05 실현) \\
⑭ & 품질 통제 기록 (§4) & 결함 1,140/1,024/116 · 밀도 0.132 > 0.13 · 감리 8건 귀속 이의 2 & 8월 SPI 0.913 정체 = 8월 결함 평탄(시험 정지) \\
⑮ & 조달·계약 변경 기록 (§5) & 라이선스 50/50 · 대기 비용 0.42 · 검수 보류 회신 8영업일 · 귀책 7+3+3$-$3 & CV 양수의 이유 — 계약 변경 로그 1번 \\
\end{tbl}
\keymsg{⑫의 어느 숫자가 다른 넷 어디에도 설명되지 않으면 그 숫자는 아직 이해되지 않은 것이다}
\note{워크북 §6 상호 관계 도식과 §0.2 20종 표. Day1의 다섯이 서로를 고쳤고 Day2의 다섯이 서로를 검산했다면 Day3의 다섯은 서로를 설명한다. 반대 방향도 말한다 — ⑫ 없이 나머지 넷은 목록에 그친다(이슈 11건이 임계경로에 무엇을 했는지, 변경 4건이 EAC를 얼마나 움직이는지, 결함 116건이 R3 개발자의 시간을 얼마나 먹는지는 성과보고에서만 한 숫자가 된다). 공통 규칙 셋을 다시 — ① 같은 사건은 같은 날짜·같은 숫자로 다섯 문서에 나타난다(R-05 실현 = IS-02 05-12 · 5월 SI SV $-$2.49 · 합의서 5호 0.42 · 귀책 7일 · S10 velocity 480) ② 착시는 지우지 않고 병기한다 ③ 모든 문서는 결정 요청으로 끝난다. 예상 소요 4분.}
\end{frame}

\begin{frame}{Day2 문서로 되돌아간 것 — 12건, 재기준선 0회}
\begin{tbl}[\scriptsize]{L{3.0cm}L{5.4cm}L{2.6cm}}
Day2 대상 & 갱신 & 결정·시점 \\ \midrule
⑥ 범위·WBS·사전 카드 5장 & v1.1 — FP 12,400 \ra{} 12,585, SI 68.20 \ra{} 69.22, 합계 145.80 \ra{} 146.82 & CCB 10-15 \ra{} BL-02 10-23 \\
⑦ RTM v1.1 & v1.2 — StRS 215 · 변경 이력 열 CR-06\til{}11 · 갱신 주기를 스프린트 리뷰 시로(감리 1번) & P1 PMO, FZ-11 \\
⑧ 일정 기준선 CPM 표 & A12 20\ra{}30일·TF 0, A09 TF 15, 2차 컷오버 03-27 추가. \textbf{A10 $-$10은 지연으로 기록}(EX-01) & P1 PM · 스폰서 10-16 \\
⑨ 원가 기준선 PV·컨틴전시 & 2026-12 +0.33 · 2027-01 +0.69 · 배정표 v1.2 · Rules of Credit 부속 2 정식 첨부 & BL-02 \\
⑩ 등록부 v1.1 · 격자 & v1.2 — 15행 재평가·R-14\til{}R-20 등재·잔여 7.26 / 원가 축 "+10.2 = 컨틴전시 전액"(CR-10) & 10-23, 정미란 서명 \\
\end{tbl}
\keymsg{기준선은 승인된 변경(BL-02)으로만 움직였다 — 지연(A10 $-$10)과 초과(밀도 0.132·클라우드 +0.6)는 편차로 남아 결정을 기다린다}
\note{워크북 §6 갱신 표 12행을 다섯 줄로 압축했다(품질 축 측정 정의·여유 소유권 문안·조정표 되돌아올 조건 해제는 노트로만). Day2 §7이 예고한 그대로 — 오늘의 측정이 어제의 기준선에 갱신을 만들었지만 재기준선(Day2 §3.5 항목 10 — CCB + 이사회, 최대 2회)은 쓰지 않았다. 그림 1-5(편차를 보이게 두기 vs 계획을 고쳐 편차 없애기)를 다시 띄워도 좋다. 학생에게 묻기 — "BL-02에 A10 $-$10을 반영하면 무엇이 달라지는가": 임계경로가 기준선 안으로 들어가 회복 계획의 근거가 사라지고, 다음 달 성과보고는 SPI(t)가 1 근처로 올라가 아무것도 재지 않게 된다. 예상 소요 3분.}
\end{frame}


\begin{frame}{Day4 예고 — 전환·종료: 오픈은 이벤트가 아니라 이관이다}
\begin{tbl}[\scriptsize]{L{0.8cm}L{4.4cm}L{5.6cm}}
 & Day4 산출물 & Day3에서 받는 입력 \\ \midrule
⑯ & 컷오버 런북 · SAC(서비스 수용 기준) & EX-01 결정(1차/2차 창) · 심각도 1·2 미해결 0 · R-14 두 번 컷오버 \\
⑰ & 운영 이관 체크리스트 · SLA/OLA/UC & 이슈 로그 미결 6건 · 감리 대응 완료 여부 · 하도급 5 \\
⑱ & 종료 보고서 · 미해결 사항 이관 목록 & BL-02 대비 최종 EAC · 변경 로그 11건 · 잔여 EMV \\
⑲ & 교훈 등록부 & 조기 신호 5개 · 착시 5개 · 감리 귀속 이의 결과 \\
⑳ & 편익 실현 원장 & 편익 축 G(118억) · B-08 리드타임 12h · 측정 책임자 실명 \\
\end{tbl}
\keymsg{Day4의 첫 질문 — "오늘 종료 보고서를 쓴다면 미해결 사항 이관 목록에 무엇이 들어가는가"}
\note{워크북 §7 "Day4로 이어지는 것" 표. Day4는 컷오버 런북(T$-$24h\til{}T+24h, go/no-go 게이트, point of no return, 롤백 트리거), ITIL Version 5 기준 서비스 전환(SAC·Early Life Support·SLA/OLA/UC 3층·Change Enablement — ITIL 4 어휘로 고정하되 현행은 v5), 편익 실현 원장 6칸(측정 책임자 실명), 종료 보고서와 교훈, 그리고 Product 트랙으로 넘어가는 접합점(이관 목록의 한 줄)을 다룬다. 오늘의 EX-01 결정이 Day4의 컷오버 창을 정하므로, Day4 아침에 "10-16에 스폰서가 무엇을 결정했는가"를 [추가 설정]으로 공개하고 시작한다. 예상 소요 3분.}
\end{frame}

\begin{frame}{사후 읽기 — 30일 안에, 그리고 Day4 전에}
\begin{itemize}
\item Day4 전(필수) — 교재 제6장 6.1 Lipke(2003) ES 원문 요약 4쪽 · 워크북 §7 표 · 자기 성과보고서 1장 재작성(SPI(t) 추가)
\item 30일 — Fleming \& Koppelman 4판 중 PMB·관리예비비 장 · Vacanti 2015 중 work item age 장 · Keil·Smith mum effect 논문 1편
\item 90일 — 자기 프로젝트의 마지막 변경 1건을 6축으로 다시 재고, CCB 의결서가 있으면 네 칸 중 비어 있는 칸을 찾는다
\item 오늘의 산출물 5종(⑪\til{}⑮)은 Day4 아침 09:00에 갤러리에 다시 붙인다 — 종료 보고서의 입력이다
\item 질문은 워크북 각 §X.7 점검 질문에 자기 답을 적어 Day4에 가져온다
\end{itemize}
\keymsg{Day3의 한 문장 — 편차는 측정으로 끝나지 않는다. 결정 요청으로 끝난다}
\note{교재 6.4 읽을 자료 표의 필수·심화(30일)·심화(90일) 구분 그대로. 서지는 리딩리스트 검증 표기를 따른다. "자기 성과보고서 재작성"은 Day4 도입에서 3명에게 발표시킨다(사전 공지). 마지막 문장은 Day1 "숫자 없는 헌장은 소설이다", Day2 "기준선은 발견되지 않는다, 협상된다"와 나란히 벽에 붙는 오늘의 명제다. 예상 소요 3분.}
\end{frame}
